\documentclass{article}

\usepackage{packages}
\usepackage{environments}
\usepackage{commands}
\usepackage{xfp} 
\usepackage{fp}

\title{Написание кода на C++}
\author{Дима Трушин}
\date{2024 --- 2025}

\AtBeginDocument{\shorthandoff{"}}


\begin{document}

\maketitle
\tableofcontents
\newpage
{
\section{State Machine или Машина Состояний (Конечный автомат)}
\colorlet{InitialStateColor}{Red!40}
\colorlet{FinalStateColor}{Blue!40}
\colorlet{EventColor}{Green!40}
\colorlet{IdleColor}{BlueGreen!40}

\subsection{Что это такое}

\subsubsection{Почти определение}

Машина состояний -- это граф с дополнительной информацией, который должен отражать текущее состояние системы и сценарии изменения этих состояний.
Чтобы в точности понять, что это такое давайте начнем с картинки:
\[
\xymatrix{
  *+=+[o][F-][F*:InitialStateColor]{\texttt{st}_\texttt{0}}\ar[r]&
  *+=+[o][F-]{\texttt{st}_\texttt{1}}\ar@/^10pt/[r]\ar[d]&
  *+=+[o][F-]{\texttt{st}_\texttt{2}}\ar@/^10pt/[l]\ar[d]&
  {}&
  *+=+[o][F-][F*:EventColor]{\texttt{e}_\texttt{0}}\ar@{}[d]|{\vdots}\\
  {}&
  *+=+[o][F-]{\texttt{st}_\texttt{3}}\ar[r]&
  *+=+[o][F-][F*:FinalStateColor]{\texttt{st}_\texttt{4}}&
  {}&
  *+=+[o][F-][F*:EventColor]{\texttt{e}_\texttt{k}}\\
}
\]
Давайте прокомментирую, что тут происходит:
\begin{enumerate}
\item У State Machine обязательно есть множество состояний.
На картинке выше -- это вершины графа.
В нашем случае это $\verb"st"_{\verb"0"}$, $\verb"st"_{\verb"1"}$,$\verb"st"_{\verb"2"}$, $\verb"st"_{\verb"3"}$ и $\verb"st"_{\verb"4"}$.

\item Среди состояний всегда должно быть в точности одно выделенное состояние, называемое Initial (начальное).
На картинке выше оно выделено $\begin{array}{|c|}\hline{\cellcolor{InitialStateColor}}\\\hline\end{array}$.

\item Среди состояний может быть выделено любое (возможно пустое) множество Final (финальных) вершин.
Их еще могут называть терминальными.
На картинке выше отмечена одна терминальная вершина  $\begin{array}{|c|}\hline{\cellcolor{FinalStateColor}}\\\hline\end{array}$.

\item Для каждой State Machine обязательно выделено некоторое множество событий (объекты из каких-то множеств), на которые State Machine может реагировать.
На картинке выше они выделены $\begin{array}{|c|}\hline{\cellcolor{EventColor}}\\\hline\end{array}$.

\item State Machine обязательно есть некоторое множество ребер, представляющие возможность перейти из одного состояния в другое.
Давайте разберем, что эти данные означают:
\[
\xymatrix{
  *+=+[o][F-]{\scriptstyle\texttt{Source}}\ar@/^10pt/[rrr]^{\raisebox{5pt}{$\scriptstyle\texttt{Trigger} \frac{\scriptstyle \texttt{[Guard]}}{\scriptstyle \texttt{Action}}$}}|(0.25){\boxed{\scriptstyle\texttt{OnOut}}}|(0.75){\boxed{\scriptstyle\texttt{OnIn}}}&{}&{}&
  *+=+[o][F-]{\scriptstyle\texttt{Target}}\\
}
\]
\begin{enumerate}
\item \verb"Source" -- состояние, из которого машина пытается перейти по ребру.

\item \verb"Target" -- состояние, в которое машина пытается перейти по ребру.

\item \verb"Trigger" -- это название сообщения, которое может запустить переход по ребру.
Обычно тут фиксируют какое-то множество из которого подходят сообщения.
В языках программирования это соответствует фиксированному типу.
В нашем случае это кто-то из $\verb"e"_{\verb"i"}$.

\item \verb"Guard" -- функция возвращающая \verb"true" или \verb"false", принимающая \verb"Source" и \verb"Trigger"%
\footnote{В оригинальном виде \texttt{Guard} должен еще принимать \texttt{Target}, но на мой взгляд это делает машину менее детерминированной, так как я не хочу смотреть в будущее при принятии решения.}
(и возможно еще некоторые данные, которые видны всем методам State Machine), которая решает допускать переход по ребру или нет.

\item \verb"Action" -- функция, которая видит:  \verb"Source", \verb"Target" и \verb"Trigger" (и возможно еще некоторые данные, которые видны всем методам State Machine), которая выполняется при переходе по ребру.
Данная функция у каждого ребра вообще говоря своя.

\item \verb"OnOut" -- функция, которая видит:%
\footnote{\label{ref::InOutNonStandard}
В оригинальной версии это действие должно еще видеть \texttt{Target} и \texttt{Trigger}, но тогда теряется логический смысл, что это действие не зависит от ребра.
Банально сообщения имеют разный тип и для обработки скажем \texttt{Trigger} нам все равно придется делать перебор по разным типам сообщений, что означает, что это скорее всего действие в рамках \texttt{Action}, которое видит все и зависит от ребра.}
\verb"Source" (и возможно еще некоторые данные, которые видны всем методам State Machine), и выполняется при выходе из \verb"Source" независимо от ребра.
Она выполняется до \verb"Action", но после \verb"Guard".

\item \verb"OnIn" -- функция, которая видит:%
\textsuperscript{\ref{ref::InOutNonStandard}}
\verb"Source" (и возможно еще некоторые данные, которые видны всем методам State Machine), и выполняется при входе в \verb"Target" независимо от ребра.
Она выполняется после \verb"Action".
\end{enumerate}

\item Когда я говорю, что какая-то функция видит скажем \verb"Source", то я подразумеваю, что \verb"Source" может хранить какие-то данные внутри.
При покидании состояния и возвращению в него у State Machine есть два вида поведения:
\begin{enumerate}
\item Мы каждый раз заходим в вершину инициализируя ее по умолчанию.

\item Мы запоминаем в каком состоянии была вершина и при возвращении в нее, инициализируем вершину сохраненными данными.
\end{enumerate}
По умолчанию мы считаем, что никакой памяти у вершин нет.
А если надо для какой-то вершины использовать память, ее надо пометить особым образом.
Я так же буду считать, что вершина \verb"Source" разрушается%
\footnote{Технически это деталь имплементации, кто где разрушается или создается.
Но эта деталь помогает думать о State Machine.
Другой вариант думать, что мы храним все вершины, но знаем, какая активная, и при входе в вершину просто сбрасываем ее содержание, если не было указано помнить историю.}
по завершению \verb"OnIn" на \verb"Target", а вершина \verb"Target" создается до запуска \verb"Guard".
Все тяжелые данные, которые не нужны для слежкой за состоянием должны быть в State Machine, а не в ее вершинах.
\end{enumerate}

\subsubsection{Функционирование}

При создании State Machine, мы обычно передаем ей данные, которые нужны ей для хранения и обработки в течение времени ее функционирования.
После создания, она находится в нерабочем состоянии, то есть ни одна вершина не является текущим состоянием.
То есть технически работающая машина состояний -- это пара из описанного выше графа с дополнительными методами и еще стрелка на вершину этого графа.%
\footnote{Во время имплементации можно как хранить текущую вершину, создавая и разрушая ее, так и хранить все вершины разом и помнить какая активная.}
И в начале эта стрелка никуда не указывает.
А потому машину состояний надо запускать явно.
При этом есть возможность:
\begin{enumerate}
\item Запустить машину с <<начальном>> состоянии.
Мы просто ставим стрелку в него.
При этом можно передать какие-то инициализирующие состояние данные.

\item Запустить машину с выбранного состояния.
При этом так же можно передать какие-то инициализирующие состояние данные.
\end{enumerate}
Как выполняется переход по ребру:
\begin{enumerate}
\item Пусть после старта State Machine находится в следующем состоянии
\[
\xymatrix{
  *+=+[o][F-][F*:EventColor]{\texttt{e}_\texttt{0}}&
  *+=+[o][F-][F*:InitialStateColor]{\texttt{st}_\texttt{0}}\ar[r]&
  *+=+[o][F-]{\texttt{st}_\texttt{1}}\ar@/^10pt/[r]\ar[d]&
  *+=+[o][F-]{\texttt{st}_\texttt{2}}\ar@/^10pt/[l]\ar[d]\\
  {}\ar@{}[ru]|(.50){}="A"|(0.8){}="B"\ar@[Plum]@{=>} "A";"B" 
  &
  {}&
  *+=+[o][F-]{\texttt{st}_\texttt{3}}\ar[r]&
  *+=+[o][F-][F*:FinalStateColor]{\texttt{st}_\texttt{4}}\\
}
\]
И прилетает событие $\verb"e"_{\verb"0"}$.
Предположим, что стрелка $\texttt{st}_\texttt{0}\to\texttt{st}_\texttt{1}$ принимает событие $\verb"e"_{\verb"0"}$.

\item Первым делом создается дефолтно инициализированное состояние $\texttt{st}_\texttt{1}$.

\item Теперь запускается проверка \verb"Guard".
Он принимает:  $\texttt{st}_\texttt{0}$,  $\texttt{st}_\texttt{1}$,  $\texttt{e}_\texttt{0}$.
Если он возвращает \verb"false", то событие отвергается и результат сообщается обработчиком события пользователю.
\[
\xymatrix{
  *+=+[o][F-][F*:EventColor]{\texttt{e}_\texttt{0}}&
  *+=+[o][F-][F*:InitialStateColor]{\texttt{st}_\texttt{0}}\ar[r]&
  *+=+[o][F-]{\texttt{st}_\texttt{1}}\ar@/^10pt/[r]\ar[d]&
  *+=+[o][F-]{\texttt{st}_\texttt{2}}\ar@/^10pt/[l]\ar[d]\\
  {}\ar@{}[ru]|(.50){}="A"|(0.8){}="B"\ar@[Plum]@{=>} "A";"B" 
  &
  {\mathmakebox[0pt][r]{\texttt{Guard(st}_\texttt{0}\texttt{, e}_\texttt{0}\texttt{)}}}&
  *+=+[o][F-]{\texttt{st}_\texttt{3}}\ar[r]&
  *+=+[o][F-][F*:FinalStateColor]{\texttt{st}_\texttt{4}}\\
}
\]

\item Если \verb"Guard" возвращает \verb"true", то событие принимается.

\item Отрабатывает \verb"OnOut" для вершины  $\texttt{st}_\texttt{0}$.
\[
\xymatrix{
  *+=+[o][F-][F*:EventColor]{\texttt{e}_\texttt{0}}&
  *+=+[o][F-][F*:InitialStateColor]{\texttt{st}_\texttt{0}}\ar[r]&
  *+=+[o][F-]{\texttt{st}_\texttt{1}}\ar@/^10pt/[r]\ar[d]&
  *+=+[o][F-]{\texttt{st}_\texttt{2}}\ar@/^10pt/[l]\ar[d]\\
  {}\ar@{}[ru]|(.50){}="A"|(0.8){}="B"\ar@[Plum]@{=>} "A";"B" 
  &
  {\mathmakebox[0pt][r]{\texttt{OnOut(st}_\texttt{0}\texttt{)}}}&
  *+=+[o][F-]{\texttt{st}_\texttt{3}}\ar[r]&
  *+=+[o][F-][F*:FinalStateColor]{\texttt{st}_\texttt{4}}\\
}
\]
\item Отрабатывает \verb"Action" для ребра.
В этот момент мы считаем, что уже покинули $\texttt{st}_\texttt{0}$ и еще не вошли в $\texttt{st}_\texttt{1}$.
\[
\xymatrix{
  *+=+[o][F-][F*:EventColor]{\texttt{e}_\texttt{0}}&
  *+=+[o][F-][F*:InitialStateColor]{\texttt{st}_\texttt{0}}\ar[r]\ar@{}[r]_(0.55){}="D"&
  *+=+[o][F-]{\texttt{st}_\texttt{1}}\ar@/^10pt/[r]\ar[d]&
  *+=+[o][F-]{\texttt{st}_\texttt{2}}\ar@/^10pt/[l]\ar[d]\\
  {}\ar@{}[ru]|(.50){}="A"|(0.8){}="B"
  \ar@[Plum]@{=>} "A" +"D" - "B"; "D"
  &
  {\mathmakebox[0pt][r]{\texttt{Action(st}_\texttt{0}\texttt{, st}_\texttt{1}\texttt{, e}_\texttt{0}\texttt{)}}}&
  *+=+[o][F-]{\texttt{st}_\texttt{3}}\ar[r]&
  *+=+[o][F-][F*:FinalStateColor]{\texttt{st}_\texttt{4}}\\
}
\]

\item Отрабатывает \verb"OnIn" для вершины  $\texttt{st}_\texttt{1}$.
\[
\xymatrix{
  *+=+[o][F-][F*:EventColor]{\texttt{e}_\texttt{0}}&
  *+=+[o][F-][F*:InitialStateColor]{\texttt{st}_\texttt{0}}\ar[r]&
  *+=+[o][F-]{\texttt{st}_\texttt{1}}\ar@/^10pt/[r]\ar[d]&
  *+=+[o][F-]{\texttt{st}_\texttt{2}}\ar@/^10pt/[l]\ar[d]\\
  {}&
  {\mathmakebox[0pt][r]{\texttt{OnIn(st}_\texttt{1}\texttt{)}}}
  \ar@{}[ru]|(.50){}="A"|(0.8){}="B"\ar@[Plum]@{=>} "A";"B" &
  *+=+[o][F-]{\texttt{st}_\texttt{3}}\ar[r]&
  *+=+[o][F-][F*:FinalStateColor]{\texttt{st}_\texttt{4}}\\
}
\]

\item Вершина $\texttt{st}_\texttt{0}$ разрушается.
И автомат приходит в состояние
\[
\xymatrix{
  {}&
  *+=+[o][F-][F*:InitialStateColor]{\texttt{st}_\texttt{0}}\ar[r]&
  *+=+[o][F-]{\texttt{st}_\texttt{1}}\ar@/^10pt/[r]\ar[d]&
  *+=+[o][F-]{\texttt{st}_\texttt{2}}\ar@/^10pt/[l]\ar[d]\\
  {} 
  &
  {}\ar@{}[ru]|(.50){}="A"|(0.8){}="B"\ar@[Plum]@{=>} "A";"B"&
  *+=+[o][F-]{\texttt{st}_\texttt{3}}\ar[r]&
  *+=+[o][F-][F*:FinalStateColor]{\texttt{st}_\texttt{4}}\\
}
\]
\end{enumerate}
Еще несколько важных замечаний по функционированию State Machine:
\begin{enumerate}
\item Если машина находится в незапущенном состоянии, она не реагирует на сообщения.

\item Если машина находится в одном из финальных состояний, она не реагирует на сообщения.

\item Если из одной вершины выходит несколько ребер, которые принимают одно и то же сообщение, то не должно быть возможно получить \verb"true" для более чем одного \verb"Guard"-а.

\item У машины состояний могут быть внутренние ребра из вершины в себя, которые не триггерят \verb"OnOut" и \verb"OnIn".
Я про эти переходы ничего не говорил, но работают они так же, просто чуть меньше действий выполняется.

\item Могут быть анонимные ребра, это такие ребра, которые не принимают никакой event в качестве триггера.
Они просто срабатывают по \verb"Guard".
Такие ребра срабатывают сразу после отработки \verb"OnIn" действия.
Если \verb"OnIn" не вызывалось (например из-за того что это внутреннее ребро), то и анонимные переходы не срабатывают.

\item Все действия, которые совершаются на ребре: \verb"Guard", \verb"OnOut", \verb"Action", \verb"OnIn" могут быть тривиальными.
Например \verb"Guard" может возвращать всегда \verb"true", а остальные действия могут просто ничего не делать.
\end{enumerate}

\subsubsection{Композиция State Machine}

Что прекрасно в машинах состояния -- их можно комбинировать разными способами, что позволяет из очень простых блоков делать достаточно хитрые механизмы для управления логикой состояния чего угодно (но нас будет интересовать именно программное обеспечение).

\paragraph{Параллельные Машины Состояния}

Если у вас есть несколько машин состояния, например $\texttt{S}_\texttt{0}$, $\texttt{S}_\texttt{1}$, то из них можно сделать машину состояний заставив работать параллельно друг с другом.
То есть вы берете тройку
\[
(\texttt{S}_\texttt{0},\;\texttt{S}_\texttt{1})
\]
И если пришло какое-то событие, то скармливаете его одновременно всем машинам состояния.
Технически надо понять, какие у этой машины состояния, какие из них начальные, какие конечные и как работают переходы.
\begin{enumerate}
\item Формально, если $S_0$ и $S_1$ -- это множества состояний этих машин соответственно, то для параллельно запущенных машин множество состояний будет декартово произведение $S_0\times S_1$.

\item Начальным состоянием будет единственная пара состоящая из начальных состояний каждой из машин.

\item Финальные состояния так же те состояния, в которых каждая машина находится в финальном состоянии.
То есть логика такая, пока хотя бы одна машина может продолжать работать, параллельная машина состояний все еще работает и может обрабатывать сообщения.

\item С ребрами будет интереснее.
Пусть нам дано по ребру из каждой машины
\[
\xymatrix{
  *+=+[o][F-]{\scriptstyle\texttt{A}_1}\ar@/^10pt/[rrr]^{\raisebox{6pt}{$\scriptstyle\texttt{E}\; \frac{\scriptstyle \texttt{[Guard}_1\texttt{]}}{\scriptstyle \texttt{Action}_1}$}}|(0.25){\boxed{\scriptstyle\texttt{Out}_1}}|(0.75){\boxed{\scriptstyle\texttt{In}_1}}&{}&{}&
  *+=+[o][F-]{\scriptstyle\texttt{B}_1}\\
}
\quad\quad\quad
\xymatrix{
  *+=+[o][F-]{\scriptstyle\texttt{A}_2}\ar@/^10pt/[rrr]^{\raisebox{6pt}{$\scriptstyle\texttt{E}\; \frac{\scriptstyle \texttt{[Guard}_2\texttt{]}}{\scriptstyle \texttt{Action}_2}$}}|(0.25){\boxed{\scriptstyle\texttt{Out}_2}}|(0.75){\boxed{\scriptstyle\texttt{In}_2}}&{}&{}&
  *+=+[o][F-]{\scriptstyle\texttt{B}_2}\\
}
\]
Здесь важно, что тип сообщения \verb"E" один и тот же для обоих ребер.
Если сообщение разные, то только одна из машин сдвинется по своему ребру, а вторая останется на месте.
Но если сообщение на обоих ребрах одинаковое, то они дают три ребра:
\[
\xymatrix@R=60pt@C=100pt{
  *++[o][F-]{\scriptstyle\texttt{A}_1\times \texttt{A}_2}
  \ar[r]^{\raisebox{6pt}{$\scriptstyle\texttt{E}\; \frac{\scriptstyle \texttt{[Guard}_1\texttt{]}}{\scriptstyle \texttt{Action}_1}$}}|(0.25){\boxed{\scriptscriptstyle \texttt{Out}_1}}|(0.75){\boxed{\scriptscriptstyle \texttt{In}_1}}
  \ar[d]_{\raisebox{6pt}{$\scriptstyle\texttt{E}\; \frac{\scriptstyle \texttt{[Guard}_2\texttt{]}}{\scriptstyle \texttt{Action}_2}$}}|(0.25){\boxed{\scriptscriptstyle \texttt{Out}_2}}|(0.75){\boxed{\scriptscriptstyle \texttt{In}_2}}
  \ar[rd]|(0.25){\boxed{\scriptscriptstyle \texttt{Out}_1,\;\texttt{Out}_2}}|(0.5){\raisebox{0pt}{$\scriptstyle\texttt{E}\; \frac{\scriptstyle \texttt{[Guard}_1\texttt{ \&\& Guard}_2\texttt{]}}{\scriptstyle \texttt{Action}_1,\;\texttt{Action}_2}$}}|(0.75){\boxed{\scriptscriptstyle \texttt{In}_1,\;\texttt{In}_2}}
  &
  *++[o][F-]{\scriptstyle\texttt{B}_1\times \texttt{A}_2}
  \\
  *++[o][F-]{\scriptstyle\texttt{A}_1\times \texttt{B}_2}
  &
  *++[o][F-]{\scriptstyle\texttt{B}_1\times \texttt{B}_2}
  \\
}
\]
То есть если сработал только $\verb"Guard"_1$, то только первая машина меняет свое состояние $\verb"A"_1$ на $\verb"B"_1$ (горизонтальная стрелка).
Если сработал только $\verb"Guard"_2$, то только вторая машина меняет свое состояние $\verb"A"_2$ на $\verb"B"_2$ (вертикальная трелка).
И только если сработали оба условия, то переход идет по диагональному ребру.

\item В предыдущем случае порядок исполнения следующией:
\begin{enumerate}
\item В начале исполняются все \verb"Guard"-ы.
После чего выбирается ребро.

\item После этого для всех вершин выполняются \verb"Out" действия.

\item После этого исполняются все \verb"Action"-ы.

\item После этого выполняются все \verb"In" действия.
\end{enumerate}


\item Стоит отметить важную вещь, что в параллельных машинах плохая идея мутировать сообщение.
Это часто запрещается, так как в целом противоречит тому, что все сайд эффекты должны быть во внутреннюю память State Machine.
При последовательном исполнении у вас машины будут видеть разные состояния сообщений, а при многопоточном исполнении, без синхронизации будут data race-ы.
\end{enumerate}

Полезный простой пример параллельных State Machine.
Вы можете стартовать с простой машины, которая включает и выключает какую-то опцию в вашей программе.
Для ее описания можно использовать машину с двумя вершинами
\[
\xymatrix{
    *+=+[o][F*:InitialStateColor][F-]{\texttt{Off}}\ar@/^10pt/[r]&
  *+=+[o][F-]{\texttt{On}}\ar@/^10pt/[l]\\
}
\]
Теперь если мы хотим добавить вторую фичу в программу, которую тоже хотим включать или выключать, то надо просто взять две такие машины работающие параллельно
\[
\xymatrix{
    {}\ar@{}[r]|(.40){}="A"|(.7){}="B"\ar@[Plum]@{=>} "A";"B"&
    *+=+[o][F*:InitialStateColor][F-]{\texttt{Off}}
    \ar@/^10pt/[r]&
    *+=+[o][F-]{\texttt{On}}
    {\save [l].[]*+[F-:<3pt>]\frm{}\restore}
    \ar@/^10pt/[l]&
    {\mathmakebox[0pt][l]{\texttt{Machine}_\texttt{1}}}\\
   {}\ar@{}[r]|(.40){}="C"|(.7){}="D"\ar@[Plum]@{=>} "C";"D"&
     *+=+[o][F*:InitialStateColor][F-]{\texttt{Off}}
    \ar@/^10pt/[r]&
    *+=+[o][F-]{\texttt{On}}
    {\save [l].[]*+[F-:<3pt>]\frm{}\restore}
    \ar@/^10pt/[l]&
    {\mathmakebox[0pt][l]{\texttt{Machine}_\texttt{2}}}\\
}
\]

\paragraph{Вложенные Машины Состояния}

Можно сделать еще более интересную вещь, а именно у машины состояний состояния сами могут быть машинами состояний.
Вот простой пример когда это может понадобиться.
Если у вас есть какая-то фича, которую можно включить и выключить, то скорее всего когда она включена у вас могут быть разные состояния.
Скажем модуль может быть включен или выключен, а когда включен он может быть в работе или на паузе.
Тогда такая ситуация хорошо описывается машиной машиной состояния с двумя узлами из предыдущего примера, где узел \verb"On" сам является машиной состояний в двумя возможными опциями \verb"On", \verb"Pause".
Например можно это изобразить так:
\[
\xymatrix{
  {\texttt{Off}}\save="C" \restore
  &
  {\phantom{\texttt{On}}}
  \ar@<3ex>@/^10pt/[r]&
  {\texttt{On}}\save="A" \restore
  &
  *+=+[o][F-]{\texttt{Play}}\ar@/^10pt/[d]\\
  {}&
  {}{}\save="D" \restore
 &
  {} \ar@<5ex>@/^10pt/[l]&
  *+=+[o][F-]{\texttt{Pause}} \save="B" \restore\ar@/^10pt/[u]
  \save"A" + <-10pt, 17pt>="AA" \restore 
   \save "B" + <15pt, -14pt>="BB"\restore
  {\save "AA"."BB"*+[F-:<3pt>]\frm{}\restore}
    \save"C" + <-10pt, 17pt>="CC" \restore 
   \save "D" + <1pt, -14pt>="DD"\restore
  {\save "CC"."DD"*+[F-:<3pt>]\frm{}\restore}
  \\
}
\]
Такая машина либо включена либо выключена, и во включенном состоянии может переключаться между тремя режимами. 
Что хорошо в этом случае -- при выключении машины, не нужно задумываться в каком она была состоянии на паузе или в работе.
Но если мы хотим, то можно сохранять состояние при выключении машины.

Давайте теперь опишем порядок действий при вложенных машинах состояния.
Рассмотрим следующий пример:
\[
\xymatrix{
  *+=+[o][F-]{\scriptstyle\texttt{A}_2}\save="A" \restore
  \ar@{}[r]^(0){\raisebox{14pt}{$\texttt{A}_1$}}
  \ar@/^10pt/[rrr]^{\raisebox{6pt}{$\scriptstyle\texttt{E}\; \frac{\scriptstyle \texttt{[Guard}_2\texttt{]}}{\scriptstyle \texttt{Action}_2}$}}|(0.25){\boxed{\scriptstyle\texttt{Out}_2}}|(0.75){\boxed{\scriptstyle\texttt{In}_2}}&{}&{}&
  *+=+[o][F-]{\scriptstyle\texttt{B}_2}\save="B" \restore
  \ar@{}@<13pt>[rrr]|(0.11){}="U"|(0.89){}="W"
  \ar@/^10pt/ "U" ; "W"^{\raisebox{6pt}{$\scriptstyle\texttt{E}\; \frac{\scriptstyle \texttt{[Guard}_1\texttt{]}}{\scriptstyle \texttt{Action}_1}$}}|(0.25){\boxed{\scriptstyle\texttt{Out}_1}}|(0.75){\boxed{\scriptstyle\texttt{In}_1}}
  &
  {}
  &
  {}
  &
  {}\save="C" \restore\ar@{}[l]_(0){\raisebox{14pt}{$\texttt{B}_1$}}
    \save"A" + <-10pt, 30pt>="AA" \restore 
   \save "B" + <4pt, -7pt>="BB"\restore
  {\save "AA"."BB"*+[F-:<3pt>]\frm{}\restore}
    \save"C" + <-10pt, 30pt>="CC" \restore 
   \save "C" + <4pt, -7pt>="DD"\restore
  {\save "CC"."DD"*+[F-:<3pt>]\frm{}\restore}
  \\
}
\]
Мы считаем, что внешняя машина находится в состоянии $\verb"A"_1$, а внутренняя в состоянии $\verb"A"_2$.
При этом обратите внимание, что мы считаем что на сообщение типа \verb"E" могут среагировать обе машины и нарисовали по одному ребру, реагирующему на данное сообщение.
В этом случае порядок действий <<снизу-вверх>>: в начале реагирует самая глубокая машина, если она не реагирует, то ее содержащая и так далее.
В нашем случае процесс исполнения будет следующий:
\begin{enumerate}
\item Вычисляется $\verb"Guard"_2$ для внутренней машины.

\item Если он дает результат \verb"true", то происходит переход на внутренней машине, а внешняя остается в $\verb"A"_1$:
\begin{enumerate}
\item Выполняются действия $\verb"Out"_2$ на выход из вершины $\verb"A"_2$.

\item Выполняется $\verb"Action"_2$.

\item Выполняются действия $\verb"In"_2$ на вход в вершину $\verb"B"_2$.
\end{enumerate}

\item Если же внутренний \verb"Guard" показал \verb"false", то управление передается внешней машине.
И проверяется $\verb"Guard"_1$.
И если он дает \verb"true", то происходит переход по ребру внешней машины.
Но обратите внимание, что вы при этом покидаете $\verb"A"_2$, а потому порядок действий будет такой:
\begin{enumerate}
\item Выполняются действия $\verb"Out"_2$ на выход из вершины $\verb"A"_2$.
% Вот это оставляет машину в некорректном состоянии.

\item Выполняются действия $\verb"Out"_1$ на выход из вершины $\verb"A"_1$.

\item Выполняется $\verb"Action"_1$.

\item Выполняются действия $\verb"In"_1$ на вход в вершину $\verb"B"_1$.
\end{enumerate}
\end{enumerate}
У внутренней машины всегда приоритет.
Например если вы жмете на клавишу \verb"Esc", то должно закрыться не все приложение, а только текущее окно.

\subsection{Зачем это нужно в программировании}

% TO DO
TO DO

\subsection{Имплементация в Compile-Time}

Теперь я хочу обсудить что и как мы хотим имплементировать.
Во-первых, я бы хотел упростить задачу и имплементировать только простую State Machine, которая не имеет другие State Machine в качестве своих состояний.
Уже для такой задачи у нас возникает несколько интересных техник, которые я бы хотел обсудить.
Во-вторых, я буду игнорировать обработку исключений.
Это не значит, что я буду использовать \verb"new/delete" всюду, но задачи \verb"OnIn" и \verb"OnOut" действий аналогичны задачам конструктора и деструктора для поддержки корректного состояния данных.
И потому между этими действиями обязана быть автоматическая связь, которую надо заворачивать в декораторы.
Так же я не стал затрагивать какие-то вопросы посвященные очереди сообщений, которая в основном нужна для вложенных State Machine, чтобы передавать информацию.
В-третьих, мне не интересно делать run-time имплементацию State Machine.
Когда вы пишете код, то у вас заранее известны компоненты вашей программы, за что они отвечают, какие ресурсы нужны для работы и по какой логике между ними переключаться.
Даже наличие динамического количества объектов можно завернуть в эту систему.
А в compile-time можно с помощью шаблонов очень сильно автоматизировать код, который еще и будет шустро работать.

\subsubsection{Описание графа для State Machine}

Начать я конечно же хочу с описания того, как я вообще хочу использовать будущую библиотеку для создания машины состояния.
Давайте начнем с конкретного примера диаграммы для State Machine:
\[
\xymatrix@R=60pt{
  *+=+[o][F*:InitialStateColor][F-]{\scriptstyle\texttt{S}_0}\ar@/^10pt/[rrr]^{\raisebox{6pt}{$\scriptstyle\texttt{E}_0\; \frac{\scriptstyle \texttt{[Guard}_0\texttt{]}}{\scriptstyle \texttt{Action}_0}$}}|(0.25){\boxed{\scriptstyle\texttt{Out}_0}}|(0.75){\boxed{\scriptstyle\texttt{In}_1}}
  &
  {}
  &
  {}
  &
  *+=+[o][F-]{\scriptstyle\texttt{S}_1}
    \ar@/^20pt/[d]^{\raisebox{0pt}{$\scriptstyle\texttt{E}_1\; \frac{\scriptstyle \texttt{[Guard}_1\texttt{]}}{\scriptstyle \texttt{Action}_1}$}}|(0.25){\boxed{\scriptstyle\texttt{Out}_1}}|(0.75){\boxed{\scriptstyle\texttt{In}_2}}
  \\
  *+=+[o][F*:FinalStateColor][F-]{\scriptstyle\texttt{S}_3}
  &
  {}
  &
  {}
  &
  *+=+[o][F-]{\scriptstyle\texttt{S}_2}
  \ar@/^10pt/[lll]^{\raisebox{-12pt}{$\scriptstyle\texttt{E}_3\; \frac{\scriptstyle \texttt{[Guard}_3\texttt{]}}{\scriptstyle \texttt{Action}_3}$}}|(0.25){\boxed{\scriptstyle\texttt{Out}_2}}|(0.75){\boxed{\scriptstyle\texttt{In}_3}}
      \ar@/^20pt/[u]^{\raisebox{0pt}{$\scriptstyle\texttt{E}_2\; \frac{\scriptstyle \texttt{[Guard}_2\texttt{]}}{\scriptstyle \texttt{Action}_2}$}}|(0.25){\boxed{\scriptstyle\texttt{Out}_2}}|(0.75){\boxed{\scriptstyle\texttt{In}_1}}
  \\
}
\]
Я буду называть это математическим представлением.

\subsubsection{Описание графа для State Machine в коде}

Давайте я опишу каким бы синтаксисом я задавал подобную машину.
Вот пример:
\begin{cppcode}
namespace detail {
using namespace SM;

class SmTable {
  constexpr auto operator()() const noexcept {
    return transition(
        *state<S<0>> >>= OnOut<0>,
        OnIn<1> >>= state<S<1>> >>= OnOut<1>,
        OnOut<2> <<= state<S<2>> <<= OnIn<2>,
        state<S<3>> <<= OnIn<3>,
        state<S<0>> + event<E<0>>[Guard<0>] / Action<0> == state<S<1>>,
        state<S<1>> + event<E<1>>[Guard<1>] / Action<1> == state<S<2>>,
        state<S<1>> <= state<S<2>> + event<E<2>>[Guard<2>] / Action<2>,
        !state<S<3>> <= state<S<2>> + event<E<3>>[Guard<3>] / Action<3>);
  }
};
} // namespace detail
\end{cppcode}
Я очень надеюсь, что происходящее угадывается, но на всякий случай подчеркну, что здесь происходит:
\begin{enumerate}
\item Прежде всего за состояния у нас отвечают отдельные классы.
Так как структура State Machine определяется в compile-time, то мы под каждое состояние заводим свой класс.
То есть состояние -- это НЕ объект класса, а именно имя класса.
В примере выше я предполагаю, что для состояний мы используем шаблон класса:
\begin{cppcode}[\nonumbers]
template<int>
struct S {};
\end{cppcode}
Слово \verb"state" -- это ключевое слово для языка построения таблицы, которое лишь проясняет назначение \verb"S<k>".

\item Под сообщения выделяются шаблоны
\begin{cppcode}[\nonumbers]
template<int>
struct E {};
\end{cppcode}
Слово \verb"event" так же является служебным и лишь поясняет роль типа \verb"E<k>".

\item Действия входа и выхода задаются шаблонными методами:
\begin{cppcode}[\nonumbers]
template<int x>
void OnIn(S<x>& state) {
  /*do something*/
}
template<int x>
void OnOut(S<x>& state) {
  /*do something*/
}
\end{cppcode}

\item Для проверки условия используются шаблонные методы
\begin{cppcode}[\nonumbers]
template<int x>
bool Guard(S<x>& state, const E<x>& event) {
  return /*something*/;
}
\end{cppcode}
Обратите внимание, что для каждого события \verb"E<x>" свой \verb"Guard", потому он знает тип события.
Порядок аргументов в методе \verb"Guard" не имеет значения, так же не имеет значения присутствуют ли оба аргумента.
Вполне возможно сделать \verb"Guard", который вообще не принимает аргументы и все время возвращает \verb"true", или принимает только сообщение или только текущее состояние.
Для обеспечения поддержки такого механизма нам понадобится диспетчеризация аргументов, которая обсуждалась в разделе~\ref{}.

\item Действия так же задаются шаблонными функциями:
\begin{cppcode}[\nonumbers]
template<int x, class = std::enable_if_t<(x != 2)>>
bool Action(S<x>& source, const E<x>& event, S<x + 1>& target) {
  return /*something*/;
}

template<int x, class = std::enable_if_t<(x == 2)>>
bool Action(S<x>& source, const E<x>& event, S<x - 1>& target) {
  return /*something*/;
}
\end{cppcode}
Не обращайте внимание на \verb"enable_if" он нужен лишь для того, чтобы для каждого индекса \verb"x" однозначно определялся метод.
Это нужно для определения только по имени какую из двух функций я имею в виду, когда пишу код вида:
\begin{cppcode}[\nonumbers]
state<S<1>>[Guard<1>] / Action<1> == state<S<2>>
\end{cppcode}
Так же как и в случае \verb"Guard" порядок аргументов у \verb"Action" не имеет значения, и более того могут использоваться не все из этих аргументов.

\item Еще одно замечание по поводу аргументов всех действий: \verb"OnIn", \verb"OnOut", \verb"Guard", \verb"Action".
Аргументы могут содержать дополнительные параметры.
В этом случае эти данные должны храниться в State Machine.
Это так называемые зависимости или данные, на которые State Machine может влиять во время своего исполнения.
Обычно машина состояний предоставляет сама внутренний механизм для хранения и инициализации подобных данных.
Эти данные так же однозначно характеризуются своим типом, то есть нельзя использовать данные одного и того же типа дважды.
Если вам нужно использовать один и тот же тип, то его надо заворачевать во wrapper с тегом идентификатором, чтобы отличит друг от друга.%
\footnote{Например подобно тому, как мы идентифицировали  type aliases в разделе~\ref{}.
Или как мы идентифицировали singleton-ы в разделе~\ref{}.}

\item Есть два вида записей:
\begin{enumerate}
\item Записи описывающие вершины с действиями \verb"OnIn" и \verb"OnOut".

\item Записи описывающие переходы между вершинами с \verb"Event", \verb"Guard" и \verb"Action".
\end{enumerate}

\item Для задания действий на вход и выход можно использовать два вида синтаксиса:
\begin{cppcode}[\nonumbers]
OnIn<1> >>= state<S<1>> >>= OnOut<1>
OnOut<1> <<= state<S<1>> <<= OnIn<1>
\end{cppcode}
При этом можно указать только одно действие из двух и другое будет тривиальным.
\begin{cppcode}[\nonumbers]
OnIn<1> >>= state<S<1>>
state<S<1>> >>= OnOut<1>
OnOut<1> <<= state<S<1>>
state<S<1>> <<= OnIn<1>
\end{cppcode}

\item Так же вершины можно отмечать как начальные (с помощью \verb"*") или как финальные (с помощью \verb"!").
Это можно делать где угодно в любых записях, можно несколько раз указать начальную вершину, главное чтобы она была одна и та же и чтобы она была.
И можно указать любое количество финальных вершин, но главное чтобы из них не было ребер.
\begin{cppcode}[\nonumbers]
*state<S<0>> >>= OnOut<0>
*state<S<0>> + event<E<0>>[Guard<0>] / Action<0> == state<S<1>>
state<S<1>> <= *state<S<0>> + event<E<0>>[Guard<0>] / Action<0>
\end{cppcode}

\item Для описания переходов используется два вида синтаксиса, с указанием назначения до или после описания перехода:%
\footnote{При этом я не трогаю оператор \texttt{operator=}, чтобы не заигрывать с семантикой языка и не получить по рогам и копытам неожиданным образом.}
\begin{cppcode}[\nonumbers]
state<S<0>> + event<E<0>>[Guard<0>] / Action<0> == state<S<1>>
state<S<1>> <= state<S<0>> + event<E<0>>[Guard<0>] / Action<0>
\end{cppcode}
При этом в таком виде можно пропускать событие, гард или действие.
Например допустимо:
\begin{cppcode}[\nonumbers]
state<S<0>> + event<E<0>>[Guard<0>] / Action<0> == state<S<1>>
state<S<0>>[Guard<0>] / Action<0> == state<S<1>>
state<S<0>> + event<E<0>> / Action<0> == state<S<1>>
state<S<0>> + event<E<0>>[Guard<0>] == state<S<1>>
state<S<0>> / Action<0> == state<S<1>>
\end{cppcode}
\end{enumerate}

\subsubsection{Добавление служебных состояний}

Теперь обратим внимание на то, что в коде я могу легко добавить действие \verb"OnIn" для начального состояния и действие \verb"OnOut" для финального:
\begin{cppcode}[\nonumbers]
*OnIn<0> >>= state<S<0>> >>= OnOut<0>
OnOut<3> <<= !state<S<3>> << OnIn<3>
\end{cppcode}
Как такое отразить в математическом представлении?
Для этого нужно ввести два служебных состояния:
\[
\xymatrix@R=60pt{
*+=+[o]{\phantom{\scriptstyle\texttt{Idle}}}&{}&
  *+=+[o][F*:InitialStateColor][F-]{\scriptstyle\texttt{I}}
  \ar[rr]|(0.6){\boxed{\scriptstyle\texttt{In}_0}}
  &{}&
  *+=+[o][F-]{\scriptstyle\texttt{S}_0}\ar@/^10pt/[rrr]^{\raisebox{6pt}{$\scriptstyle\texttt{E}_0\; \frac{\scriptstyle \texttt{[Guard}_0\texttt{]}}{\scriptstyle \texttt{Action}_0}$}}|(0.25){\boxed{\scriptstyle\texttt{Out}_0}}|(0.75){\boxed{\scriptstyle\texttt{In}_1}}
  &
  {}
  &
  {}
  &
  *+=+[o][F-]{\scriptstyle\texttt{S}_1}
    \ar@/^20pt/[d]^{\raisebox{0pt}{$\scriptstyle\texttt{E}_1\; \frac{\scriptstyle \texttt{[Guard}_1\texttt{]}}{\scriptstyle \texttt{Action}_1}$}}|(0.25){\boxed{\scriptstyle\texttt{Out}_1}}|(0.75){\boxed{\scriptstyle\texttt{In}_2}}
  \\
  {}&{}&
  *+=+[o][F*:FinalStateColor][F-]{\scriptstyle\texttt{F}}
  &{}&
  *+=+[o][F-]{\scriptstyle\texttt{S}_3}
  \ar[ll]|(0.4){\boxed{\scriptstyle\texttt{Out}_3}}
  &
  {}
  &
  {}
  &
  *+=+[o][F-]{\scriptstyle\texttt{S}_2}
  \ar@/^10pt/[lll]^{\raisebox{-12pt}{$\scriptstyle\texttt{E}_3\; \frac{\scriptstyle \texttt{[Guard}_3\texttt{]}}{\scriptstyle \texttt{Action}_3}$}}|(0.25){\boxed{\scriptstyle\texttt{Out}_2}}|(0.75){\boxed{\scriptstyle\texttt{In}_3}}
      \ar@/^20pt/[u]^{\raisebox{0pt}{$\scriptstyle\texttt{E}_2\; \frac{\scriptstyle \texttt{[Guard}_2\texttt{]}}{\scriptstyle \texttt{Action}_2}$}}|(0.25){\boxed{\scriptstyle\texttt{Out}_2}}|(0.75){\boxed{\scriptstyle\texttt{In}_1}}
  \\
}
\]
Для начального состояния \verb"S<0>" добавляем вершину \verb"I", которая становится новым начальным состоянием и делаем анонимный%
\footnote{То есть этот переход тригерится автоматически без какого либо события сразу после входа в состояние \texttt{I}.}
переход из \verb"I" в \verb"S<0>" на которое подвешивается \verb"OnIn" для вершины \verb"S<0>".
Аналогично для \verb"S<3>" добавляем новую вершину \verb"F", которая становится новым финальным состоянием, и делаем анонимный переход из \verb"S<3>" в \verb"F", на который подвешиваем \verb"OnOut" для \verb"S<3>".

Есть еще одно служебное состояние \verb"Idle", которое нам понадобится.
Это состояние нужно для того, чтобы описать не запущенную State Machine.
Оно так же служебное и вводится неявно.
Его обычно не добавляют в граф, либо можно нарисовать его отдельной изолированной вершиной без действия на в ход и выход.
\[
\xymatrix@R=60pt{
  *+=+[o][F*:IdleColor][F-]{\scriptstyle\texttt{Idle}}&{}&
  *+=+[o][F*:InitialStateColor][F-]{\scriptstyle\texttt{I}}
  \ar[rr]|(0.6){\boxed{\scriptstyle\texttt{In}_0}}
  &{}&
  *+=+[o][F-]{\scriptstyle\texttt{S}_0}\ar@/^10pt/[rrr]^{\raisebox{6pt}{$\scriptstyle\texttt{E}_0\; \frac{\scriptstyle \texttt{[Guard}_0\texttt{]}}{\scriptstyle \texttt{Action}_0}$}}|(0.25){\boxed{\scriptstyle\texttt{Out}_0}}|(0.75){\boxed{\scriptstyle\texttt{In}_1}}
  &
  {}
  &
  {}
  &
  *+=+[o][F-]{\scriptstyle\texttt{S}_1}
    \ar@/^20pt/[d]^{\raisebox{0pt}{$\scriptstyle\texttt{E}_1\; \frac{\scriptstyle \texttt{[Guard}_1\texttt{]}}{\scriptstyle \texttt{Action}_1}$}}|(0.25){\boxed{\scriptstyle\texttt{Out}_1}}|(0.75){\boxed{\scriptstyle\texttt{In}_2}}
  \\
  {}&{}&
  *+=+[o][F*:FinalStateColor][F-]{\scriptstyle\texttt{F}}
  &{}&
  *+=+[o][F-]{\scriptstyle\texttt{S}_3}
  \ar[ll]|(0.4){\boxed{\scriptstyle\texttt{Out}_3}}
  &
  {}
  &
  {}
  &
  *+=+[o][F-]{\scriptstyle\texttt{S}_2}
  \ar@/^10pt/[lll]^{\raisebox{-12pt}{$\scriptstyle\texttt{E}_3\; \frac{\scriptstyle \texttt{[Guard}_3\texttt{]}}{\scriptstyle \texttt{Action}_3}$}}|(0.25){\boxed{\scriptstyle\texttt{Out}_2}}|(0.75){\boxed{\scriptstyle\texttt{In}_3}}
      \ar@/^20pt/[u]^{\raisebox{0pt}{$\scriptstyle\texttt{E}_2\; \frac{\scriptstyle \texttt{[Guard}_2\texttt{]}}{\scriptstyle \texttt{Action}_2}$}}|(0.25){\boxed{\scriptstyle\texttt{Out}_2}}|(0.75){\boxed{\scriptstyle\texttt{In}_1}}
  \\
}
\]

\subsubsection{Создание и использование машины состояний}

Теперь как я вижу создание самой машины.
Предположим, что мы описали таблицу переходов для машины как это сделано выше в коде классом \verb"SmTable".
Теперь можно создать машину так:
\begin{cppcode}
namespace detail {
using namespace SM;

class SmTable {
  constexpr auto operator()() const noexcept {
    return transition(/*description of the transition table*/);
  }
};
} // namespace detail

int main() {

  using Machine = SM::StateMachine<detail::SmTable>;

  Machine s;         // starts at Idle state
  s.process(E<0>{}); // ignores machine is not started

  s.start();         // jumps to I and automatically passes to S<0>

  s.process(E<0>{}); // passes to S<1>
  s.process(E<0>{}); // ignores and stays at S<1>
  s.process(E<1>{}); // passes to S<2>
  s.process(E<2>{}); // passes to S<1>
  s.process(E<1>{}); // passes to S<2>
  s.process(E<3>{}); // passes to S<3> and automatically passes to F
  return 0;
}
\end{cppcode}
Я использую служебное пространство имен \verb"detail" для того, чтобы спрятать \verb"SmTable", потому что \verb"using namespace" нельзя использовать внутри класса, только внутри другого пространства имен.
Это нужно, чтобы при описании таблицы переходов писать просто \verb"state<S<0>>" или \verb"event<E<0>>" вместо \verb"SM::state<S<0>>" и \verb"SM::event<E<0>>" соответственно.
А так как я не хочу раскрывать служебные данные из \verb"SM" (пространства имен, где лежит публичный код для использования State Machine), то приходится пользоваться дополнительным пространством имен \verb"detail".

\subsection{Общий план имплементации}

Обратите внимание на то, как в коде выше у нас процесс создания State Machine разбит на два шага:
\begin{enumerate}
\item Описание таблицы переходов.
\begin{cppcode}[\nonumbers]
namespace detail {
using namespace SM;

class SmTable {
  constexpr auto operator()() const noexcept {
    return transition(/*description of the transition table*/);
  }
};
} // namespace detail
\end{cppcode}

\item Создание State Machine по таблице переходов.
\begin{cppcode}[\nonumbers]
using Machine = SM::StateMachine<detail::SmTable>;
\end{cppcode}
\end{enumerate}
Так можно разделить код на две части:
\begin{enumerate}
\item Публичная часть (Front), которая предоставляет пользователю механизмы для описания таблицы переходов.
Например, служебные слова \verb"state" и \verb"event", а так же операторы \verb"[]", \verb"+" и \verb"/" и т.д.
Так же эта часть кода делает проверки на корректность данных для задания таблицы.
В данном случае делаются не содержательные проверки на устройство State Machine, а лишь синтаксические проверки, что пользователь использует только доступные механизмы для описания и использует их корректно.

\item Приватная часть (Back), которая используя таблицу созданную пользователем непосредственно создает класс самой State Machine.
Эта часть кода делает проверки корректности полученной State Machine.
Например, что машина имеет в точности одно начальное состояние, что анонимные переходы не имеют циклов и прочее.
\end{enumerate}
Ниже я буду придерживаться такого же разделения на FrontEnd и BackEnd.
Это позволит локализовать контекст решения задач.
Для создания Front мы по сути хотим создать свой язык внутри C++ для описания таблиц.
Это так называемый Domain Specific Language (DSL).
Описание таблицы

\begin{cppcode}
namespace detail {
using namespace SM;

class SmTable {
  constexpr auto operator()() const noexcept {
    return transition(
        *state<S<0>> >>= OnOut<0>,
        OnIn<1> >>= state<S<1>> >>= OnOut<1>,
        OnOut<2> <<= state<S<2>> <<= OnIn<2>,
        state<S<3>> <<= OnIn<3>,
        state<S<0>> + event<E<0>>[Guard<0>] / Action<0> == state<S<1>>,
        state<S<1>> + event<E<1>>[Guard<1>] / Action<1> == state<S<2>>,
        state<S<1>> <= state<S<2>> + event<E<2>>[Guard<2>] / Action<2>,
        !state<S<3>> <= state<S<2>> + event<E<3>>[Guard<3>] / Action<3>);
  }
};
} // namespace detail
\end{cppcode}
является примером DSL, где мы с помощью операторов описываем таблицу переходов в очень понятном виде.%
\footnote{Я считаю это удивительным совпадением.
Если бы еще была возможность использовать оператор \texttt{=>} вместо \texttt{==}, то было вы вообще идеально.}
Я хочу построить таблицу перехода (пока не важно как она задается) в compile-time.
Есть два подхода:
\begin{enumerate}
\item Использовать шаблонные классы.

\item Использовать шаблонные \verb"constexpr" переменные (спасибо C++17).
\end{enumerate}
Я буду пользоваться вторым подходом.
Его главное преимущество заключается в том, что с переменными в отличие от типов можно использовать операторы, которые напоминают синтаксис для задания State Machine.
При этом нас не будут интересовать значения переменных, а информация будет содержаться в типе переменной.

Сама таблица переходов у меня будет состоять из следующих записей
\begin{cppcode}[\nonumbers]
(OnEntry, State, OnExit)
(Source, Event, Guard, Action, Target)
(State, Type)
\end{cppcode}
Первая запись хранит имя состояния и действия на вход и выход для этого состояния.
Вторая запись отвечает за переход.
А третья запись для каждой вершины знает ее тип (например что это начальная или финальная вершина или \verb"Idle" состояние и прочее).
Состояния, события и тип состояния -- это просто классы, потому их можно хранить в списке типов, которые мы обсуждали в разделе~\ref{}.
А функции \verb"OnEntry", \verb"OnExit", \verb"Guard", \verb"Action" можно спрятать внутрь шаблонного типа можно используя вспомогательный шаблон для хранения типизированного элемента:
\begin{cppcode}[\nonumbers]
template<class T, T v>
struct Value_t {
  using value_type = T;
  static constexpr T value = v;
};
\end{cppcode}
Тогда можно писать
\begin{cppcode}[\nonumbers]
using Guard = Value_t<decltype(&guard_function), &guard_function>;
\end{cppcode}
Это по сути тот же самый подход, который я использовал для хранения списка элементов одного и того же типа, только у нас внутри самого шаблона есть служебные записи для более удобного восстановления аргументов.
Например мне не надо будет писать метод доступа аналогичный \verb"EL::At" для получения элемента, а можно будет просто написать \verb"Guard::value" и все.
Еще я вместо хранения пар \verb"(State, Type)" буду отдельно хранить имя начальной вершины и список имен финальных вершин.

Идея имплементации машины состояний такая.
Мы будем хранить внутри только текущую вершину для текущего состояния.
Так как все типы заранее известны (включая служебные вершины), то можно использовать \verb"std::variant" и Visitor паттерн для работы с вершинами.
Вот как может выглядеть шаблон для самой State Machine:
\begin{cppcode}
// TEdges        -- List of transitions of the form (Source, Event, Guard, Action, Target)
// EntryExit     -- List of entries of the form (OnEntry, State, OnExit)
// InitialStates -- List of initial states
// TFinalStates  -- List of final states
// TEvents       -- List of all events
template<class Edges, class EntryExit, class InitialStates, class TFinalStates, class TEvents>
class SM<Edges, EntryExit, InitStates, TFinalStates, TEvents> {
public:
  using Edges = /*...*/;        // Add edges for pseudostates
  using States =  /*...*/;      // Extract states from Edges and EntryExit lists
  using Idle = Back::IdleState; // The Idle state, MUST not be used by the user

private:
  using StatesExt = TL::PushFront<Idle, I, F, States>;
  // Extend the States with Idle, I, F states
  using AState = TL::Convert<StatesExt, std::variant>;
  // AState = std::variant<Idle, I, F, S<0>, S<1>, S<2>, S<3>>;

public:
  using Initial = TL::Front<InitStates>; // There must be only one
  using FinalStates = TFinalStates;
  using Events = TEvents;

  template<class Event>
  bool process(Event&& event); // the main processing function

  template<class State>
  bool isAt() const; // check the state

  void reset();      // move to idle state
  bool start();      // start from I state

  bool isComplete() const;
  bool isStarted() const;
  bool isRunning() const;

  template<class Event>
  static constexpr bool acceptEvent() noexcept;
  template<class State>
  static constexpr bool isValidState() noexcept;
  template<class State>
  static constexpr bool isFinalState() noexcept;
  template<class State>
  static constexpr bool isInitialState() noexcept;

private:
  template<class Event>
  bool processOnce(Event&& event);
  template<class State>
  bool forceMoveTo(State&& target)

  template<class State>
  bool processAnonymousOnce();
  bool processAnonymousOnce();
  template<class State>
  bool processAnonymous();
  bool processAnonymous();

  AState state_ = Idle{}; // Holds the state of the machine. Idle is the initial one.
};
\end{cppcode}
Так как мы имеем дело с \verb"std::variant", то основной метод для работы со состоянием -- это создание объектов визитеров для определения того, кто лежит внутри и выполнения соответствующих действий уже зная текущий тип состояния.
Типичная рутина для обработки сообщения \verb"process" будет выглядеть как-то так:
\begin{enumerate}
\item Проверить \verb"isRunning" (машина в не \verb"Idle" и не в \verb"F" состоянии).
Если машина не запущена, то вернуть \verb"false", если запущена идем далее.

\item Проверить обрабатывается сообщение машиной или нет (содержится ли в списке \verb"Events").
Если нет, то возвращаем \verb"false".
Если да, то идем дальше.

\item Выполнить обработку события \verb"Event" с помощью визитера.
Если обработчик вернул \verb"false", значит перехода по ребру не произошло (например не сработал ни один \verb"Guard") и мы тут же возвращаем \verb"false".
Если обработчик вернул \verb"true" идем дальше.

\item Теперь надо обработать анонимные переходы.
Это делаем в цикле пока \verb"processAnonyousOnce" возвращает \verb"true".
После чего возвращаем \verb"true" независимо от того, были ли сделаны анонимные переходы.
\end{enumerate}
При написании обработчика события можно на самом деле сделать три разные имплементации:
\begin{enumerate}
\item Описанная выше имплементация, когда на каждый единичный переход по ребру мы запускаем визитера.
То есть на каждое ребро у нас будет отдельный \verb"switch" по всем возможным состояниям из \verb"std::variant".

\item Если мы один раз зашли внутрь \verb"std::variant" и знаем таблицу переходов, то мы точно знаем куда мы можем перейти.
А потому можно не выходить из визитера, а рекурсивно пройтись по всем ребрам во время компиляции.
Это избавляет нас от лишнего \verb"switch" на каждое ребро, но добавляет рекурсию.

\item Можно использовать очередь служебных сообщений для управления State Machine и идти не в цикле по анонимным перходам, а в цикле пока очередь событий не будет пустой.
\end{enumerate}
Я предпочел первый вариант, так как он самый простой и понятный.
К нему так же не сложно добавить очередь.
Рекурсивное решение так же не трудно написать, но оно требует чуть больше интересных решений на уровне шаблонов.

\subsection{FrontEnd Создание DSL для генерации таблицы}

Давайте в начале опишем, что мы хотим сделать и как.
Я не хочу опускаться до формального уровня описания грамматики языка.%
\footnote{Не знаю на сколько это к добру.}
Наша задача уметь задавать таблицы перехода для State Machine в следующем виде
\begin{cppcode}
class SmTable {
  constexpr auto operator()() const noexcept {
    entry1, entry2, entry3,...., entryn
  }
};
\end{cppcode}
Здесь \verb"entryk" -- допустимые записи в таблице.
Теперь нужно объяснить какие записи мы допускаем.

\subsubsection{Записи с одной вершиной}

Первый тип записей поддерживает описание информации связанной с отдельно взятой вершиной.
В эту информацию входят:
\begin{enumerate}
\item Класс отвечающий за состояние.

\item Класс отвечающий за статус состояния (начальная вершина, финальная вершина или обычная%
\footnote{Возможны и другие состояния в будущем.}
).

\item Действие \verb"OnEntry" выполняемое на вход в состояние.

\item Действие \verb"OnExit" выполняемое на выходе из состояния.
\end{enumerate}
Суть последних двух действий в том, что они работают в паре как конструктор/деструктор и отвечают за корректность состояния данных для работы внутри текущего состояния.

Теперь опишем синтаксис для пользователя, где приведем все возможные варианты:
\begin{center}
\begin{tabular}{lll}
{
\begin{minipage}[\baselineskip]{5.7cm}
\begin{cppcode}[\nonumbers]
state<S>
state<S> >>= OnExit
OnExit <<= state<S>
state<S> <<= OnEntry
OnEntry >>= state<S>
OnEntry >>= state<S> >>= OnExit
OnExit <<= state<S> <<= OnEntry
\end{cppcode}
\end{minipage}
}&{
\begin{minipage}[\baselineskip]{5.7cm}
\begin{cppcode}[\nonumbers]
*state<S>
*state<S> >>= OnExit
OnExit <<= *state<S>
*state<S> <<= OnEntry
OnEntry >>= *state<S>
OnEntry >>= *state<S> >>= OnExit
OnExit <<= *state<S> <<= OnEntry
\end{cppcode}
\end{minipage}
}&{
\begin{minipage}[\baselineskip]{5.7cm}
\begin{cppcode}[\nonumbers]
!state<S>
!state<S> >>= OnExit
OnExit <<= !state<S>
!state<S> <<= OnEntry
OnEntry >>= !state<S>
OnEntry >>= !state<S> >>= OnExit
OnExit <<= !state<S> <<= OnEntry
\end{cppcode}
\end{minipage}
}\\
\end{tabular}
\end{center}
Опишем что есть что.
\begin{enumerate}
\item Символ \verb"S" означает любой класс.
Данные классы должны однозначно идентифицировать состояния машины состояний.

\item \verb"state" -- служебная шаблонная переменная доступная пользователю в публичной части интерфейса библиотеки для создания State Machine.

\item \verb"OnEntry" -- метод возвращающий \verb"void", который может принимать \verb"S" (текущее состояние).
Принимать он может его как по значению, так и по ссылке на усмотрение пользователя, то есть доступные варианты: \verb"S", \verb"S&", \verb"const S&", \verb"S&&".%
\footnote{Возможно \texttt{S\&\&} -- не самый разумный вариант, ибо вы поглощаете состояние, над которым исполняете какую-то инициализацию.
То есть это состояние исчезнет для внешнего наблюдателя.
Так что по-хорошему такую сигнатуру надо запретить, но я забью на это для простоты.}
Метод может не принимать текущее состояние вовсе.
Так же метод может принимать любые другие классы, главное, чтобы они все имели разный тип с точностью до снятия ссылки и \verb"const/volatile" квалификаторов верхнего уровня.%
\footnote{Два типа будут считаться одинаковыми, если \texttt{std::is\_same\_v<std::remove\_cvref\_t<T>, std::remove\_cvref\_t<U>{}>}  возвращает \texttt{true}.}
Все дополнительные аргументы должны содержаться внутри State Machine и передаваться с помощью механизма диспетчеризации аргументов для State Machine.%
\footnote{Я позволю себе считать, что у меня нет дополнительных внутренних данных, опять же для простоты имплементации.}

\item Оператор \verb"*" означает, что состояние помечено как начальное.
Оператор \verb"!" означает, что состояние помечено как финальное.
Отсутствие любого из операторов выше помечает состояние как обычное.
Знак \verb"*" или \verb"!" достаточно поставить у вершины только один раз.
Если хотя бы раз вершина помечена таким знаком, она становится начальной или финальной соответственно.
Можно пометить таким знаком одну и ту же вершину несколько раз.
Единственное ограничение -- нельзя помечать одну и ту же вершину обоими знаками в разных записях.

\item Записи
\begin{cppcode}[\nonumbers]
state<S> >>= OnExit
OnExit <<= state<S>
\end{cppcode}
устанавливают действие на выходе из состояния.

\item Записи
\begin{cppcode}[\nonumbers]
state<S> <<= OnEntry
OnEntry >>= state<S>
\end{cppcode}
устанавливают действие на входе в состояние.

\item Записи 
\begin{cppcode}[\nonumbers]
OnExit <<= state<S> <<= OnEntry
OnEntry >>= state<S> >>= OnExit
\end{cppcode}
одновременно устанавливают оба действия.

\item Допускается для одной и той же вершины несколько раз установить входное (или выходное) действие, но нужно, чтобы во всех случаях это была одна и та же функция.
Можно вовсе не устанавливать никаких действий, тогда система должна их определить как тривиальные.
Это служебная функция, которая не принимает ничего и не делает ничего.
\begin{cppcode}[\nonumbers]
constexpr bool always_true() noexcept {
  return true;
}
\end{cppcode}
\end{enumerate}

\subsubsection{Имплементация операторов для одного состояния}

Прежде всего нам понадобится три вспомогательных класса для обозначения статуса состояния:
\begin{cppcode}[\nonumbers]
struct Regular;
struct Initial;
struct Final;

using StatusList = TL::List<Regular, Initial, Final>;
\end{cppcode}
Добавляем список всех поддерживаемых состояний на будущее.
Теперь нам нужен класс для хранения информации о вершине.
\begin{cppcode}[\nonumbers]
template<class S, class Status>
struct Status {};

template<class S>
constexpr State<S, Regular> state = {};
\end{cppcode}
Таким образом код
\begin{cppcode}[\nonumbers]
state<S>
\end{cppcode}
создает переменную с типом 
\begin{cppcode}[\nonumbers]
State<S, Regular>
\end{cppcode}
Чтобы можно было сменить статус вершины, нужно определить два оператора:
\begin{cppcode}[\nonumbers]
template<class S, class Status>
struct Status {
  constexpr State<S, Initial> operator*() const noexcept {
    return State<S, Initial>{};
  }
  constexpr State<S, Final> operator!() const noexcept {
    return State<S, Final>{};
  }
};
\end{cppcode}
Теперь надо заставить работать строчи
\begin{cppcode}[\nonumbers]
state<S> >>= OnExit
OnEntry >>= state<S>
\end{cppcode}
Кроме того, раз мы установили действие на выход (или вход), то нет смысла его устанавливать еще раз.
Например мы не хотим иметь возможность написать так:
\begin{cppcode}[\nonumbers]
OnExit1 <<= state<S> >>= OnExit2
OnEntry1 >>= state<S> <<= OnEntry2
\end{cppcode}
А это значит, что у выражения
\begin{cppcode}[\nonumbers]
state<S> >>= OnExit
\end{cppcode}
не должно быть оператора
\begin{cppcode}[\nonumbers]
?? operator<<=(Ex on_exit, state<S>)
\end{cppcode}
Но может быть оператор
\begin{cppcode}[\nonumbers]
?? operator<<=(state<S>, En on_entry)
\end{cppcode}
А для выражения
\begin{cppcode}[\nonumbers]
OnEntry >>= state<S>
\end{cppcode}
наоборот.
А потому эти операторы будут возвращать разные типы.
Я буду их называть по тому, какая информация в них хранится.
\begin{cppcode}[\nonumbers]
template<class S, class St, class Ex>
struct StateExit {
  Ex on_exit;
};

template<class S, class St, class En>
struct StateEntry {
  En on_entry;
};
\end{cppcode}
Теперь я могу задать указанные операторы так:
\begin{cppcode}[\nonumbers]
template<class S, class St>
struct State {
  ...
  template<class Ex, class = std::enable_if_t<!areSame2<State<S, St>, std::remove_cv_t<Ex>>>>
  constexpr StateExit<S, St, Ex> operator>>=(Ex on_exit) const noexcept {
    return {on_exit};
  }
  ...
};

template<class S, class St, class En,
         class = std::enable_if_t<!areSame2<State<S, St>, std::remove_cv_t<En>>>>
constexpr StateEntry<S, St, En> operator>>=(En on_entry, State<S, St>) noexcept {
  return {on_entry};
}
\end{cppcode}
Так как все эти операции будут выполняться во время компиляции, я не переживаю и передаю все аргументы по значению, а не по ссылке или как-то еще.
Я бы мог воспользоваться \verb"consteval", чтобы сделать исполнение только в compile-time принудительно, но что-то я не стал так делать, а надо бы.
Тут же можно добавить операторы \verb"<<=".
\begin{cppcode}[\nonumbers]
template<class S, class St>
struct State {
  ...
  template<class En, class = std::enable_if_t<!areSame2<State<S, St>, std::remove_cv_t<En>>>>
  constexpr StateEntry<S, St, En> operator<<=(En on_entry) const noexcept {
    return on_entry >>= *this;
  }
  ...
};

template<class S, class St, class Ex,
         class = std::enable_if_t<!areSame2<State<S, St>, std::remove_cv_t<Ex>>>>
constexpr StateExit<S, St, Ex> operator<<=(Ex on_exit, State<S, St> s) noexcept {
  return s >>= on_exit;
}
\end{cppcode}
Обратите внимание, что я не пишу имплементацию заново, а свожу к уже определенному оператору.
Теперь добавим полные данные
\begin{cppcode}[\nonumbers]
template<class S, class St, class En, class Ex>
struct StateEntryExit {
  En on_entry;
  Ex on_exit;
};
\end{cppcode}
Теперь если у нас есть уже метод входа в состояние, то добавление методов выхода выглядит так:
\begin{cppcode}[\nonumbers]
template<class S, class St, class En>
struct StateEntry {
  template<class Ex>
  constexpr StateEntryExit<S, St, En, Ex>
  operator>>=(Ex on_exit) const noexcept {
    return {on_entry, on_exit};
  }
  En on_entry;
};

template<class S, class St, class En, class Ex>
constexpr StateEntryExit<S, St, En, Ex>
operator<<=(Ex on_exit, StateEntry<S, St, En> s) noexcept {
  return s >>= on_exit;
}
\end{cppcode}
Аналогично делается добавление входного действия, когда уже есть выходное.
Я не буду писать уже этот код тут.
Что важно, все созданные типы не конвертируются друг к другу.




\subsubsection{Записи для переходов таблицы}

Второй тип записей описывает ребра в машине состояний.
В информацию о ребрах входит следующая информация:
\begin{enumerate}
\item Класс отвечающий за состояние \verb"Source" для начала перехода.

\item Класс отвечающий за статус состояния \verb"Source".

\item Класс \verb"Event" отвечающий за тип события.
Объекты такого типа являются триггерами для данного перехода.
Если ребро является анонимным переходом (нет события), то мы будем использовать \verb"void".

\item \verb"Guard" -- функция возвращающая \verb"bool".
Она может принимать текущее состояние и сообщение (как по значению, так и по ссылке).
Так же может принимать любые другие данные.
Ограничение на аргументы такое же как и у \verb"OnEntry" и \verb"OnExit": все аргументы должны быть разного типа (после снятия ссылок и \verb"const/reference" квалификаторов самого верхнего уровня).

\item \verb"Action" -- функция возвращающая \verb"void".
В качестве аргументов может принимать \verb"Source" -- состояние откуда происходит переход, \verb"Target" -- состояние куда происходи переход, \verb"Event" -- соответствующее сообщение.
Все аргументы могут приниматься как по значению, так и по любому виду ссылок.
Какие-то из этих аргументов могут отсутствовать.
Порядок аргументов не важен.
Могут присутствовать еще другие данные, которые хранятся в State Machine.
Требование на аргументы такое же как и выше: все аргументы должны быть разного типа (после снятия ссылок и \verb"const/reference" квалификаторов самого верхнего уровня).

\item Класс \verb"Target" отвечающий за состояние куда происходит переход.

\item Класс отвечающий за статус нового состояния \verb"Target".
\end{enumerate}

Теперь опишем возможные для пользователя операции.
Здесь я проигнорирую операторы \verb"*" и \verb"!".
Предполагается, что их можно добавить к любой записи вида \verb"state<S>".
Доступны постфиксная и префиксная записи для перехода из вершины \verb"S1" в \verb"S2":
\begin{cppcode}
// postfix Source ... == Target
state<S1> + event<Event> [Guard] / Action == state<S2>
state<S1> + event<Event> [Guard] == state<S2>
state<S1> + event<Event> / Action == state<S2>
state<S1> + event<Event>  == state<S2>
state<S1> [Guard] / Action == state<S2>
state<S1> [Guard]  == state<S2>
state<S1> / Action == state<S2>
// prefix Target <= Source ...
state<S2> <= state<S1> + event<Event> [Guard] / Action
state<S2> <= state<S1> + event<Event> [Guard]
state<S2> <= state<S1> + event<Event> / Action
state<S2> <= state<S1> + event<Event>
state<S2> <= state<S1> [Guard] / Action
state<S2> <= state<S1> [Guard]
state<S2> <= state<S1> / Action
\end{cppcode}
Опишем, что здесь означает что и какие есть требования к компонентам:
\begin{enumerate}
\item \verb"S1" обозначает тип класса идентифицирующего состояние из которого происходит переход.

\item \verb"S2" означает тип класса идентифицирующего состояние в которое происходит переход.

\item \verb"state" -- служебная шаблонная переменная для указания роли типа \verb"S" в выражении \verb"state<S>".
Значит, что \verb"S" идентифицирует состояние.

\item \verb"Event" -- тип сообщений объекты которого являются триггерами перехода.

\item \verb"event" -- служебная шаблонная переменная для указания роли типа \verb"Event" в выражении \verb"event<Event>".
Значит, что \verb"Event" идентифицирует тип сообщения.

\item \verb"Guard" --  функция возвращающая \verb"bool".
Может принимать в качестве аргументов состояние \verb"S1" и сообщение \verb"Event".
Так же может принимать и другие данные.
Требования на аргументы такие же, как и при описании \verb"OnEntry" и \verb"OnExity" в разделе выше.
\item Если \verb"Guard" не указан, то используется служебная функция всегда возвращающая \verb"true"
\begin{cppcode}[\nonumbers]
constexpr bool always_true() noexcept {
  return true;
}
\end{cppcode}

\item \verb"Action" -- функция возвращающая \verb"void".
Может принимать в качестве аргументов состояние \verb"S1", сообщение \verb"Event" и новое состояние \verb"S2".
Может принимать и другие данные.
Требования к аргументам такие же как и у \verb"Guard".
\item Если \verb"Action" не указан, то используется служебная функция возвращающая \verb"void" и ничего не делающая
\begin{cppcode}[\nonumbers]
constexpr void do_nothing() noexcept {
}
\end{cppcode}

\item В указанных выражениях допустимо, чтобы \verb"S1" и \verb"S2" был один и тот же тип.
Однако, в этом случае обязательно должно быть указано либо сообщение, либо гард, то есть возможны следующие опции:
\begin{cppcode}
// postfix Source ... == Target
state<S> + event<Event> [Guard] / Action == state<S>
state<S> + event<Event> [Guard] == state<S>
state<S> + event<Event> / Action == state<S>
state<S> + event<Event>  == state<S>
state<S> [Guard] / Action == state<S>
state<S> [Guard]  == state<S>
// prefix Target <= Source ...
state<S> <= state<S> + event<Event> [Guard] / Action
state<S> <= state<S> + event<Event> [Guard]
state<S> <= state<S> + event<Event> / Action
state<S> <= state<S> + event<Event>
state<S> <= state<S> [Guard] / Action
state<S> <= state<S> [Guard]
\end{cppcode}
Последний переход может показаться странным, потому что мы проверяем \verb"Guard" и без каких либо действий перезаходим в текущее состояние.
Однако, перезаход в текущее состояние тригерит два действия по умолчанию: \verb"OnExit" и \verb"OnEntry".
Такое действие по сути сбрасывает текущее состояние в начальное положение.
Если же допустить переход вида
\begin{cppcode}[\nonumbers]
// postfix Source ... == Target
state<S> / Action  == state<S>
// prefix Target <= Source ...
state<S> <= state<S> / Action
\end{cppcode}
то он приведет к бесконечному исполнению.
Без события и гарда, такое действие будет крутиться бесконечно.
\end{enumerate}
Теперь я хочу добавить еще один вид переходов -- это внутренние переходы для состояния.
Это такой переход из состояния в себя, который не тригерит \verb"OnExit" и \verb"OnExit" действия.
Внутренний переход задается без указания состояния назначения так:
\begin{cppcode}[\nonumbers]
state<S1> + event<Event> [Guard] / Action
state<S1> + event<Event> / Action
state<S1> [Guard] / Action
\end{cppcode}
Обратите внимание, что во всех случаях присутствует \verb"Action".
Действительно для внутреннего состояния нет смысла поддерживать переходы без действий, потому что в этом случае \verb"OnExit" и \verb"OnEntry" не тригерятся, \verb"Action" ничего не делает, а это значит что данные никак не меняются.
Но и состояние State Machine тоже не меняется.
А это означает, что нет смысла их поддерживать, это просто лишнее исполнение.
Кроме того, обязательно должен присутствовать либо \verb"Guard" либо сообщение (либо оба сразу).
Потому что анонимные переходы без \verb"Guard" сами в себя будут крутиться в бесконечной рекурсии.

\subsubsection{Имплементация операторов кодирующих переходы}

Здесь я хочу начать с того, чтобы имплементировать все операторы, чтобы скомпилировалась строчка
\begin{cppcode}[\nonumbers]
state<S1> + event<E>[Guard] / Action == state<S2>
\end{cppcode}
В данном случае происходит вычисление операторов в следующем порядке
\begin{cppcode}[\nonumbers]
            event<E>[Guard]
            event<E>[Guard] / Action
state<S1> + event<E>[Guard] / Action
state<S1> + event<E>[Guard] / Action == state<S2>
\end{cppcode}
Давайте заведем дополнительно класс для сохранения  информации о сообщении.
\begin{cppcode}[\nonumbers]
template<class E>
struct Event {
};
\end{cppcode}
Так же нам понадобятся
\begin{cppcode}[\nonumbers]
template<class E, class G>
struct EventGuard {
  G guard;
};

template<class E, class G, class A>
struct EventGuardAction {
  G guard;
  A action;
};
\end{cppcode}
Теперь добавим оператор \verb"operator[]" так:
\begin{cppcode}[\nonumbers]
template<class E>
struct Event {
  template<class G>
  constexpr EventGuard<E, G> operator[](G guard) const noexcept {
    return {guard};
  }
};
\end{cppcode}
Теперь оператор добавления действия
\begin{cppcode}[\nonumbers]
template<class E, class G>
struct EventGuard {
  template<class A>
  constexpr EventGuardAction<E, G, A> operator/(A action) const noexcept {
      return {guard, action};
  }
  G guard;
};
\end{cppcode}
Давайте сразу обработаем случай, когда нет \verb"Guard".
В этом случае надо уметь обработать выражение
\begin{cppcode}[\nonumbers]
event<E> / Action
\end{cppcode}
В этом случае надо определить оператор \verb"operator/" сразу для \verb"Event".
\begin{cppcode}[\nonumbers]

template<class E>
struct Event {
  template<class A>
  constexpr EventGuardAction<E, Core::TGuard, A>
  operator/(A action) const noexcept {
    return {always_true, action};
  }
};
\end{cppcode}
Тут сразу возвращается тройка вместе с тривиальным \verb"Guard"-ом.
Теперь нужен оператор сложения между состоянием и тем, что могло быть построено выше.
То есть надо уметь складывать выражения вида
\begin{cppcode}[\nonumbers]
state<S> + event<E>[Guard] / Action
state<S> + event<E>[Guard]
state<S> + event<E>
\end{cppcode}
Вначале определим классы для хранения результата:
\begin{cppcode}[\nonumbers]
template<class S, class St, class E>
struct StateEvent {
};
template<class S, class St, class E, class G>
struct StateEventGuard {
  G guard;
};
template<class S, class St, class E, class G, class A>
struct InternalEdge {
  G guard;
  A action;
};
\end{cppcode}
А теперь добавляем три оператора в класс \verb"State":
\begin{cppcode}[\nonumbers]
template<class S, class St>
struct State {
  ...
  template<class E, class G, class A>
  constexpr InternalEdge<S, St, E, G, A>
  operator+(EventGuardAction<E, G, A> edg) const noexcept {
    return {edg.guard, edg.action};
  }

  template<class E, class G>
  constexpr StateEventGuard<S, St, E, G>
  operator+(EventGuard<E, G> eg) const noexcept {
    return {eg.guard};
  }

  template<class E>
  constexpr StateEvent<S, St, E> operator+(Event<E>) const noexcept {
    return StateEvent<S, St, E>{};
  }
  ...
};
\end{cppcode}
Теперь нужно научиться обрабатывать ситуации следующего вида:
\begin{cppcode}
state<S1>[Guard] / Action
state<S1> Action
\end{cppcode}
Для этого добавляем оператор вида:
\begin{cppcode}[\nonumbers]
template<class S, class St>
struct State {
  template<class G>
  constexpr StateEventGuard<S, St, void, G> operator[](G guard) const noexcept {
    return {guard};
  }
  template<class A>
  constexpr InternalEdge<S, St, void, Core::TGuard, A>
  operator/(A action) const noexcept {
    return {always_true, action};
  }
};
\end{cppcode}
После этого мы умеем строить любые выражение в левой части от оператора \verb"operator==" и настало время определить сам оператор.
Давайте перечислим все что может и что не может стоять слева от этого оператора.
Слева выражения, а справа их типы:
\begin{cppcode}[\nonumbers]
// Ok
state<S> + event<E>[Guard] / Action // -> InternalEdge<S, St, E, G, A>
state<S> + event<E>[Guard]          // -> StateEventGuard<S, St, E, G>
state<S> + event<E> / Action        // -> InternalEdge<S, St, E, G, A>
state<S> + event<E>                 // -> StateEvent<S, St, E>
state<S>[Guard] / Action            // -> InternalEdge<S, St, E, G, A> 
state<S>[Guard]                     // -> StateEventGuard<S, St, E, G>
state<S> / Action                   // -> InternalEdge<S, St, E, G, A> 
// Almost OK
state<S>                            // -> State<S, St>
// Bad
event<E>[Guard] / Action            // -> EventGuardAction<E, G, A>
event<E>[Guard]                     // -> EventGuard<E, G>
event<E> Action                     // -> EventGuardAction<E, G, A>
event<E>                            // -> Event<E>
\end{cppcode}
Второй аргумент для оператора \verb"operator==" всегда должен быть \verb"State<S, St>", а потому мы лишь должны перебрать все левые аргументы.
Обратите внимание, что \verb"State<S, St>" почти разрешен как левый аргумент для \verb"operator==".
А именно, он разрешен если ребро ведет в другое состояние.
Как обычно в начале надо завести класс для хранения результата
\begin{cppcode}[\nonumbers]
template<class S, class St, class E, class G, class A, class T, class TSt>
struct Edge {
  G guard;
  A action;
};
\end{cppcode}
Теперь, имплементацию оператора \verb"operator==" сделаем следующим образом.
Добавим этот оператор во все классы выше, но для тех классов, для которых этот оператор должен быть не определен, поставим \verb"static_assert" с сообщением, что нельзя задать переход с такими условиями.
Теперь имплементации:
\begin{cppcode}[\nonumbers]
template<class S, class St>
struct State {
  ...
  template<class S1, class St1>
  constexpr Edge<S, St, void, Guard, Action, S1, St1> operator==(State<S1, St1>) const noexcept {
    static_assert(!areSame2<std::remove_cv_t<S>, std::remove_cv_t<S1>>,
                  "Cannot define anonymous transitions from a state to itself!");
    return {always_true, do_nothing};
  }
  ...
};
\end{cppcode}
Здесь \verb"Guard" -- это тип стандартного гарда \verb"always_true", а \verb"Action" -- это тип стандартного пустого действия \verb"do_nothing".
Теперь для разных видов состояний с дополнительной информацией:
\begin{cppcode}[\nonumbers]
template<class S, class St, class E, class G>
struct StateEventGuard {
  template<class S1, class St1>
  constexpr Edge<S, St, E, G, Action, S1, St1>
  operator==(State<S1, St1>) const noexcept {
    return {guard, do_nothing};
  }
  G guard;
};

template<class S, class St, class E>
struct StateEvent {
  template<class S1, class St1>
  constexpr Edge<S, St, E, Guard, Action, S1, St1> operator==(State<S1, St1>) const noexcept {
    return {always_true, do_nothing};
  }
};

template<class S, class St, class E, class G, class A>
struct InternalEdge  {
  template<class S1, class St1>
  constexpr Edge<S, St, E, G, A, S1, St1> operator==(State<S1, St1>) const noexcept {
    return {guard, action};
  }
  G guard;
  A action;
};
\end{cppcode}
И теперь для <<плохих>> случаев:
\begin{cppcode}[\nonumbers]
template<class E>
struct Event {
  template<class S, class St>
  constexpr Edge<void, void, E, Guard, Action, S, St> operator==(State<S, St>) const noexcept {
    static_assert(False<E>, "Cannot define a transition without source!");
    return {always_true, do_nothing};
  }
};

template<class E, class G>
struct EventGuard {
  template<class S, class St>
  constexpr Edge<void, void, E, G, Action, S, St> operator==(State<S, St>) const noexcept {
    static_assert(False<E>, "Cannot define a transition without source!");
    return {guard, do_nothing};
  }
};

template<class E, class G, class A>
struct EventGuardAction  {
  template<class T, class ST>
  constexpr Edge<void, void, E, G, A, T, ST>  operator==(State<T, ST>) const noexcept {
    static_assert(False<E>, "Cannot define a transition without source!");
    return {guard, action};
  }
};
\end{cppcode}
Еще надо не забыть имплементировать оператор \verb"<=", но его можно целиком засунуть внутрь класса \verb"State" и перегрузить через \verb"operator==".

\subsubsection{Имплементация внутренних переходов}

В этом случае мы хотим допустить следующие записи в таблице переходов (мы это обсуждали выше):%
\footnote{Данная таблица объясняет почему я назвал получающийся класс не \texttt{StateEventGuardAction} а \texttt{InternalEdge}.}
\begin{cppcode}[\nonumbers]
state<S1> + event<Event> [Guard] / Action // -> InternalEdge<S, St, E, G, A>
state<S1> + event<Event> / Action         // -> InternalEdge<S, St, E, G, A>
state<S1> [Guard] / Action                // -> InternalEdge<S, St, E, G, A>
\end{cppcode}
При этом остаются выражения, которые нам надо запретить следующего вида
\begin{cppcode}[\nonumbers]
state<S1> + event<Event>[Guard] // -> StateEventGuard<S, St, E, G>
state<S1> + event<Event>        // -> StateEvent<S, St, E>
state<S1>[Guard]                // -> StateEventGuard<S, St, E, G>
state<S1> / Action              // -> InternalEdge<S, St, E, G, A>
state<S1>                       // -> State<S, St>

event<Event>[Guard] / Action    // -> EventGuardAction<E, G, A>
event<Event>[Guard]             // -> EventGuard<E, G>
event<Event> / Action           // -> EventGuardAction<E, G, A>
event<Event>                    // -> Event<E>
\end{cppcode}
Обратите внимание, что мы можем легко запретить по типу все случаи кроме одного:
\begin{cppcode}[\nonumbers]
state<S1> / Action              // -> InternalEdge<S, St, E, G, A>
\end{cppcode}
Его можно отличить лишь исследовав само ребро для внутреннего перехода.
В этом случае событие \verb"E" будет \verb"void" и гард для этого перехода будет тривиальный \verb"always_true".
Так как возвращаемые значения не подставляются больше ни в какой оператор, то отловить эти случаи так же как мы обрабатывали \verb"operator==" уже не получится.
А потому я перекину эту работу на следующую часть кода, которая будет конвертировать данные в списки типов.

\subsubsection{Конвертация в список типов}
\label{section::SMFrontToBackConvert}

Теперь давайте скажем, что у нас получилось.
Для начала надо сказать, что из себя представляет функция \verb"transition".
Эта функция просто возвращает tuple от своих аргументов.
\begin{cppcode}[\nonumbers]
template<class... Data>
constexpr Tuple<Data...> transition(Data... data) noexcept {
  using T = Tuple<Data...>;
  return T(data...);
}
\end{cppcode}
Вот тут я использую свою кастомную \verb"constexpr" версию \verb"Tuple" элементы которой можно использовать в качестве шаблонных параметров (смотри раздел~\ref{section::Tuple}).
После вызова этой функции, все методы \verb"OnEntry", \verb"OnExit", \verb"Guard" и \verb"Action" хранятся в полях элементов, а не в самих типах.
А для удобной работы с будущей таблицей на уровне BackEnd-а, нужно сделать все единообразным и запаковать данные из \verb"Tuple" внутрь типов.

\paragraph{Что я хочу}

Давайте я продемонстрирую на примере, что я хочу сделать и как.
Вот например у меня есть запись в таблице
\begin{cppcode}[\nonumbers]
state<S1> + event<E>[guard] / action == state<S2>
\end{cppcode}
где \verb"guard" и \verb"action" -- некоторые функции.
Тогда
\begin{cppcode}[\nonumbers]
transition(state<S1> + event<E>[guard] / action == state<S2>)
\end{cppcode}
возвращает объект
\begin{cppcode}[\nonumbers]
Edge<S1, St1, E, G, A, S2, St2> t;
assert(t.guard == guard);
assert(t.action== action);
\end{cppcode}
Здесь \verb"G" -- тип указателя на функцию \verb"guard", а \verb"A" -- тип указателя на функцию \verb"action".
И мы видим, что типы для этих действий хранятся в виде шаблонных параметров, но сами функции как указатели в run-time параметрах.
Я же хочу перетащить их внутрь шаблонных параметров так.
То есть по элементу \verb"t" я хочу построить следующий тип:
\begin{cppcode}[\nonumbers]
TL::List<S1, St1, E, Value_t<G, guard>, Value_t<A, action>, S2, St2>
\end{cppcode}

Теперь рассмотрим другой вид записи
\begin{cppcode}[\nonumbers]
on_entry >>= state<S> >>= on_exit
\end{cppcode}
Где \verb"on_entry" и \verb"on_exit" -- некоторые функции.
Теперь
\begin{cppcode}[\nonumbers]
transition(on_entry >>= state<S> >>= on_exit)
\end{cppcode}
возвращает объект
возвращает объект
\begin{cppcode}[\nonumbers]
StateEntryExit<S, St, En, Ex> t;
assert(t.on_entry == entry);
assert(t.on_exit == on_exit);
\end{cppcode}
Я хочу заменить элемент \verb"t" на тип следующего вида:
\begin{cppcode}[\nonumbers]
TL::List<Value_t<En, on_entry>, S, St, Value_t<Ex, on_exit>>;
\end{cppcode}
где \verb"En" -- тип функции \verb"on_entry", а \verb"Ex" -- тип функции \verb"on_exit".
На выходе у меня будет список содержащий только следующие записи:
\begin{cppcode}[\nonumbers]
TL::List<S1, St1, E, Value_t<G, guard>, Value_t<A, action>, S2, St2>
TL::List<Value_t<En, on_entry>, S, St, Value_t<Ex, on_exit>>
\end{cppcode}
Дальше я его разделю на два списка \verb"Edges" и \verb"EnExList" (для первых и вторых записей соответственно).
После чего отдельно составлю списки:
\begin{enumerate}
\item Всех вершин отмеченных хотя бы раз как начальная.

\item всех вершин отмеченных хотя бы раз как конечная.
\end{enumerate}
И уже после этого вырежу статусы вершин из данных и у меня будут два списка: \verb"Edge" будет содержать записи вида:
\begin{cppcode}[\nonumbers]
TL::List<S1, E, Value_t<G, guard>, Value_t<A, action>, S2>
\end{cppcode}
и \verb"EnExList" записи вида
\begin{cppcode}[\nonumbers]
TL::List<Value_t<En, on_entry>, S, Value_t<Ex, on_exit>>
\end{cppcode}

\paragraph{Операции над \texttt{EnExList}}

Для упрощение работы с записями вида
\begin{cppcode}[\nonumbers]
TL::List<Value_t<En, on_entry>, S, St, Value_t<Ex, on_exit>>
\end{cppcode}
Я определю вспомогательные команды, которые сделают код более читаемым:
\begin{cppcode}[\nonumbers]
namespace StateOp {
template<class Node>
using State = TL::At<Node, 1>;

template<class Node>
using Status = TL::At<Node, 2>;

template<class Node>
using OnEntryW = TL::At<Node, 0>;

template<class Node>
using OnEntryType = typename OnEntryW<Node>::value_type;

template<class Node>
constexpr auto OnEntry = OnEntryW<Node>::value;

template<class Node>
using OnExitW = TL::At<Node, 3>;

template<class Node>
using OnExitType = typename OnExitW<Node>::value_type;

template<class Node>
constexpr auto OnExit = OnExitW<Node>::value;

template<class T>
struct Bind {
  template<class Node>
  using is_t = Logic::areSame2_t<T, Status<Node>>;
  template<class Node>
  static constexpr bool is = is_t<Node>::value;

  template<class Node>
  using isSame_t = Logic::areSame2_t<T, State<Node>>;
  template<class Node>
  static constexpr bool isSame = isSame_t<Node>::value;
};
} // namespace StateOp
\end{cppcode}
Здесь из пояснений разве что только требует шаблон \verb"Bind" с его внутренними шаблонами.
Мне нужны следующие два предиката:
\begin{enumerate}
\item Статус \verb"St" вершины в данной записи есть \verb"T".

\item Состояние \verb"S" в данной записи есть \verb"T".
\end{enumerate}
Эти предикаты зависят от параметра \verb"T".
И вот \verb"Bind" нужен как раз для того, чтобы запомнить этот параметр.
Вот как это работает.
Пусть у нас есть запись
\begin{cppcode}[\nonumbers]
using Node = TL::List<Value_t<En, on_entry>, S, St, Value_t<Ex, on_exit>>;
\end{cppcode}
Тогда
\begin{cppcode}[\nonumbers]
static constexpr bool res = Bind<S>::template isSame<Node>;
// res is true
static constexpr bool res1 = Bind<S1>::template isSame<Node>;
// res is false if S is not the same type as S1
\end{cppcode}
И аналогично для статуса
\begin{cppcode}[\nonumbers]
static constexpr bool res = Bind<St>::template is<Node>;
// res is true
static constexpr bool res1 = Bind<St1>::template is<Node>;
// res is false if St is not the same type as St1
\end{cppcode}
Самый большой недостаток такого синтаксиса в том, что приходится писать назойливое слово \verb"template" после \verb"::".
Но тут ничего не поделаешь.
Мы этот код будем использовать во внутреннем коде и такой синтаксис хотя бы не утечет на сторону пользователя.

\paragraph{Проверки тривиальности действий}

Теперь нам нужно уметь понимать действия \verb"on_entry" и \verb"on_exit" тривиальные или нет.
То есть совпадают они с нашей функцией \verb"do_nothing" или нет.
Для этого надо ввести вспомогательные шаблоны проверки условия равенства.
Однако, тут мы столкнемся с интересной особенностью сравнения указателей на функцию.
А именно, нельзя сравнивать указатели разной сигнатуры.
Попытка сравнить два указателя даже если они отличаются флагом \verb"noexcept" влечет ошибку компиляции, а потому для сравнения указателей на функцию придется писать свой метод сравнения.
Сначала сравниваем по сигнатуре, если они разные, то \verb"false".
А если они одинаковые, то тогда можно сравнить сами указатели.
Вот как это должно выглядеть.

Прежде всего вот как выглядит код для служебных методов:
\begin{cppcode}
constexpr bool always_true() noexcept {
  return true;
}
using TGuard = decltype(&always_true);

constexpr void do_nothing() noexcept {
}
using TAction = decltype(&do_nothing);

using NoAction = Value_t<TAction, do_nothing>;
using NoGuard = Value_t<TGuard, always_true>;
\end{cppcode}
А вот так выглядят методы сравнения указателей на функции
\begin{cppcode}
template<class A, class B>
constexpr std::enable_if_t<!areSame2<A, B>, bool> areSameFunctions(A,
                                                                   B) noexcept {
  return false;
}

template<class A>
constexpr bool areSameFunctions(A f, A g) noexcept {
  return f == g;
}
\end{cppcode}
Тогда проверку, что \verb"on_entry" действие тривиально имплементируется так:
\begin{cppcode}
namespace detail {
template<class Node, bool isCorrect = /*checks*/>
struct isOnEntryTrivialImpl {
  /*error messages*/
};

template<class Node>
struct isOnEntryTrivialImpl<Node, true> : Bool_t<areSameFunctions(OnEntry<Node>, do_nothing)> {};
} // namespace detail

template<class Node>
struct isOnEntryTrivial_t : detail::isOnEntryTrivialImpl<Node> {};

template<class Node>
constexpr bool isOnEntryTrivial = isOnEntryTrivial_t<Node>::value;
\end{cppcode}

Еще мне понадобятся проверки, что запись является записью про \verb"on_entry" и \verb"on_exit" действия.
К таким записям относятся элементы \verb"Tuple" следующего вида:
\begin{cppcode}
State<S, St>
StateEntry<S, St, En>
StateExit<S, St, Ex>
StateEntryExit<S, St, En, Ex>
\end{cppcode}
Я буду считать, что у нас есть следующие шаблоны для проверки:
\begin{cppcode}
template<class T>
struct isState_t : /*implementation*/ {};
template<class T>
constexpr bool isState = isState_t<T>::value;

template<class T>
struct isStateEntry_t : /*implementation*/ {};
template<class T>
constexpr bool isStateEntry = isStateEntry_t<T>::value;

template<class T>
struct isStateExit_t : /*implementation*/ {};
template<class T>
constexpr bool isStateExit = isStateExit_t<T>::value;

template<class T>
struct isStateEntryExit_t : /*implementation*/ {};
template<class T>
constexpr bool isStateEntryExit = isStateEntryExit_t<T>::value;

template<class T>
struct isAnyState_t : Or_w<isState_t<T>, isStateExit_t<T>,
                           isStateEntry_t<T>, isStateEntryExit_t<T>> {};

template<class T>
constexpr bool isAnyState = isAnyState_t<T>::value;
\end{cppcode}

\paragraph{Операции над ребрами}

Для упрощения работы с ребрами, то есть записями вида:
\begin{cppcode}[\nonumbers]
TL::List<S1, St1, E, Value_t<G, guard>, Value_t<A, action>, S2, St2>
\end{cppcode}
Так же определим вспомогательные операции.
\begin{cppcode}

namespace EdgeOp {
template<class Edge>
using Source = TL::At<Edge, 0>;

template<class Edge>
using SourceStatus = TL::At<Edge, 1>;

template<class Edge>
using Event = TL::At<Edge, 2>;

template<class Edge>
using GuardW = TL::At<Edge, 3>;

template<class Edge>
using GuardType = typename GuardW<Edge>::value_type;

template<class Edge>
constexpr auto Guard = GuardW<Edge>::value;

template<class Edge>
using ActionW = TL::At<Edge, 4>;

template<class Edge>
using ActionType = typename ActionW<Edge>::value_type;

template<class Edge>
constexpr auto Action = ActionW<Edge>::value;

template<class Edge>
using Target =
    IfElse<areSame2<TL::At<Edge, 5>, void>, Source<Edge>, TL::At<Edge, 5>>;

template<class Edge>
using TargetStatus = TL::At<Edge, 6>;

template<class Edge>
struct isInternal_t : areSame2_t<TL::At<Edge, 5>, void> {};

template<class Edge>
constexpr bool isInternal = isInternal_t<Edge>::value;

template<class Edge>
constexpr bool isAnonymous = areSame2<Event<Edge>, void>;

template<class Edge>
struct isAnonymous_t : Bool_t<isAnonymous<Edge>> {};

template<class T>
struct Bind {
  template<class Edge>
  using isSource_t = Logic::areSame2_t<T, std::remove_cv_t<SourceStatus<Edge>>>;
  template<class Edge>
  static constexpr bool isSource = isSource_t<Edge>::value;

  template<class Edge>
  using isSourceSame_t = Logic::areSame2_t<T, std::remove_cv_t<Source<Edge>>>;
  template<class Edge>
  static constexpr bool isSourceSame = isSourceSame_t<Edge>::value;

  template<class Edge>
  using isTarget_t = Logic::areSame2_t<T, std::remove_cv_t<TargetStatus<Edge>>>;
  template<class Edge>
  static constexpr bool isTarget = isTarget_t<Edge>::value;

  template<class Edge>
  using isTargetSame_t = Logic::areSame2_t<T, std::remove_cv_t<Target<Edge>>>;
  template<class Edge>
  static constexpr bool isTargetSame = isTargetSame_t<Edge>::value;
};
} // namespace EdgeOp
\end{cppcode}
Тут все более или менее тоже самое, только для ребра у нас по два шаблона внутри \verb"Bind", так как у ребра две вершины.

Еще нам понадобятся шаблоны для проверки является ли запись ребром.
При этом у нас может быть два вида ребер с записями следующих типов:
\begin{cppcode}[\nonumbers]
Edge<S1, St1, E, G, A, S2, St2>
InternalEdge<S1, St1, E, G, A>
\end{cppcode}
И я буду считать, что у нас есть шаблоны для проверки
\begin{cppcode}
template<class T>
struct isEdge_t : /*implementation*/ {};
template<class T>
constexpr bool isEdge = isEdge_t<T>::value;

template<class T>
struct isInternalEdge_t : /*implementation*/ {};
template<class T>
constexpr bool isInternalEdge = isInternalEdge_t<T>::value;

template<class T>
struct isAnyEdge_t : Or_w<isEdge_t<T>, isInternalEdge_t<T>> {};
template<class T>
constexpr bool isAnyEdge = isAnyEdge_t<T>::value;
\end{cppcode}
С помощью последних двух методов я смогу спокойно отбирать из \verb"Tuple" записей соответствующих любому из двух видов ребер.

\paragraph{Перепаковка \texttt{EnExList} записей}

Здесь мы будем перерабатывать запись из \verb"Tuple" в тип.
Публичный интерфейс
\begin{cppcode}
template<auto st>
struct PackStateInType_t
    : detail::PackStateInTypeImpl<std::remove_cv_t<decltype(st)>, st> {};

template<auto st>
using PackStateInType = typename PackStateInType_t<st>::type;
\end{cppcode}
Общая имплементация используется для сообщения об ошибках
\begin{cppcode}
namespace detail {
template<class S, auto st>
struct PackStateInTypeImpl {
  static_assert(False<S>, "The argument of PackStateInType MUST be a state from DSL: "
                "State, StateExit, StateEntry, StateEntryExit!");
};
} // namespace detail
\end{cppcode}
Перепаковка происходит в специализациях.
Так как у нас приходит $4$ типа записей сюда, то надо написать $4$ специализации:
\begin{cppcode}
namespace detail {
template<class S, class St, State<S, St> st>
struct PackStateInTypeImpl<State<S, St>, st> {
  using type = TL::List<NoAction, S, St, NoAction>;
  using value_type = type;
  static constexpr value_type value = type{};
};

template<class S, class St, class Ex, StateExit<S, St, Ex> st>
struct PackStateInTypeImpl<StateExit<S, St, Ex>, st> {
  using type = TL::List<NoAction, S, St, Value_t<Ex, st.on_exit>>;
  using value_type = type;
  static constexpr value_type value = type{};
};

template<class F, class St, class En, StateEntry<S, St, En> st>
struct PackStateInTypeImpl<StateEntry<S, St, En>, st> {
  using type = TL::List<Value_t<En, st.on_entry>, S, St, NoAction>;
  using value_type = type;
  static constexpr value_type value = type{};
};

template<class S, class St, class En, class Ex, StateEntryExit<S, St, En, Ex> st>
struct PackStateInTypeImpl<StateEntryExit<S, St, En, Ex>, st> {
  using type = TL::List<Value_t<En, st.on_entry>, S, St, Value_t<Ex, st.on_exit>>;
  using value_type = type;
  static constexpr value_type value = type{};
};
} // namespace detail
\end{cppcode}

\paragraph{Перепаковка ребер}

Перепаковка ребер делается по той же схеме, только опций всего две, а не четыре.
Мы также хотим перепаковать запись из \verb"Tuple" в шаблон с параметрами.
Публичный интерфейс будет такой:
\begin{cppcode}
template<auto edge>
struct PackEdgeInType_t
    : detail::PackEdgeInTypeImpl<std::remove_cv_t<decltype(edge)>, edge> {};

template<auto edge>
using PackEdgeInType = typename PackEdgeInType_t<edge>::type;
\end{cppcode}
Как обычно общая имплементация используется для сообщения об ошибках:
\begin{cppcode}
namespace detail {
template<class E, auto edge>
struct PackEdgeInTypeImpl {
  static_assert(False<E>, "The argument of PackInType MUST be and Edge or InternalEdge!");
};
} // namespace detail
\end{cppcode}
И теперь две специализации для каждого вида ребер:
\begin{cppcode}
template<class S1, class St1, class E, class G, class A, class S2, class St2,
         Edge<S1, St1, E, G, A, S2, St2> edge>
struct PackEdgeInTypeImpl<Edge<S1, St1, E, G, A, S2, St2>, edge> {
  using type = TL::List<S1, St1, E, Value_t<G, edge.guard>, Value_t<A, edge.action>, S2, St2>;
  using value_type = type;
  static constexpr value_type value = type{};
};

template<class S, class St, class E, class G, class A, InternalEdge<S, St, E, G, A> edge>
struct PackEdgeInTypeImpl<InternalEdge<S, St, E, G, A>, edge> {
  using type = TL::List<S, St, E, Value_t<G, edge.guard>, Value_t<A, edge.action>, void, void>;
  using value_type = type;
  static constexpr value_type value = type{};
};
\end{cppcode}
Обратите внимание, что именно в этом месте мы выделяем внутренние переходы с помощью \verb"void" в качестве состояния назначения ребра.

\paragraph{Перепаковка данных для BackEnd}

Теперь как устроена перепаковка таблицы.
У нас дан класс для генерации таблицы
\begin{cppcode}
struct SmTable {
  constexpr auto operator()() const noexcept {
    return transition(/*table description*/);
  }
};
\end{cppcode}
Теперь мы хотим взять \verb"SmTable" и вернуть два списка.
Первый содержит записи вида:
\begin{cppcode}
TL::List<Value_t<En, on_entry>, S, St, Value_t<Ex, on_exit>>;
\end{cppcode}
Второй содержит записи вида:
\begin{cppcode}
TL::List<S1, St1, E, Value_t<G, guard>, Value_t<A, action>, S2, St2>
\end{cppcode}
Для этого заведем специальный шаблон:
\begin{cppcode}
template<class Table>
struct ExtractTransitionData_t {
  static constexpr auto data = Table{}();

  template<auto x>
  struct IsState : DSL::isAnyState_t<std::remove_cv_t<decltype(x)>> {};
  template<auto x>
  struct IsEdge : DSL::isAnyEdge_t<std::remove_cv_t<decltype(x)>> {};

  static constexpr auto s_data = Filter<IsState, data>;
  static constexpr auto e_data = Filter<IsEdge, data>;

  using StateList = TL::NoDuplicates<
      TL::ConvertToList<TL::RemoveCVRec<Apply_t<PackStateInType_t, s_data>>>>;
  using EdgeList = TL::NoDuplicates<
      TL::ConvertToList<TL::RemoveCVRec<Apply_t<PackEdgeInType_t, e_data>>>>;
  using type = TL::List<StateList, EdgeList>;
};

template<class Table>
using ExtractTransitionData = typename ExtractTransitionData_t<Table>::type;
\end{cppcode}
Давайте прокомментирую код
\begin{enumerate}
\item В начале мы создаем \verb"Tuple" содержащий всю информацию вызывая оператор \verb"operator()" на объекте \verb"Table{}".

\item Чтобы вытащить из \verb"Tuple" другую \verb"Tuple" состоящую только из записей одного вида нужен предикат.
У нас уже есть предикат от типа элемента, но для \verb"Tuple" нужен предикат от самого элемента.
Тут мы определяем такие предикаты \verb"IsState" и \verb"IsEdge".

\item Далее мы создаем два отдельных \verb"Tuple" содержащих только записи отдельного вида.

\item После чего идут две распаковки.
Давайте разберемся с первой из них, вторая работает аналогично:
\begin{cppcode}
TL::NoDuplicates<
    TL::ConvertToList<
        TL::RemoveCVRec<
            Apply_t<PackStateInType_t, s_data>
        >
    >
>;
\end{cppcode}
Пусть у нас \verb"s_data" выглядит так:
\begin{cppcode}
Tuple<StateEntryExit<S1, St1, En1, Ex1>, StateEntryExit<S2, St2, En2, Ex2>> s_data;
\end{cppcode}
Это значит, что у нас две записи.
При этом внутри каждой записи мы храним действия \verb"on_entry" и \verb"on_exit".
После применения \verb"Apply_t" (см.~\ref{}) получим тип (\verb"Apply" возвращает новый \verb"Tuple", а \verb"Apply_t" возвращает тип объекта):
\begin{cppcode}
Tuple<TL::List<Value_t<En1, on_entry1>, S1, St1, Value_t<Ex1, on_exit1>>,
      TL::List<Value_t<En2, on_entry2>, S2, St2, Value_t<Ex2, on_exit2>>>
\end{cppcode}
Дальше мы рекурсивно убираем все \verb"const/volatile" квалификаторы.
Потом преобразуем внешний \verb"Tuple" к \verb"TL::List".
И под конец убираем дублирующиеся записи.

\item В самом конце возвращается пара из двух списков
\begin{cppcode}
TL::List<EnExList, EdgeList>
\end{cppcode}
\end{enumerate}
Для доступа к последним листам я сделаю методы доступа:
\begin{cppcode}
template<class Data>
using EnExList = TL::At<Data, 0>;
template<class Data>
using EdgeList = TL::At<Data, 1>;
\end{cppcode}

Здесь я хочу добавить еще небольшой кусок кода для проверки.
Дело в том, что сейчас мы проигнорировали в таблице все записи, которые не подходят под предикаты \verb"IsState" и \verb"IsEdge".
А они могли быть, если пользователь вбил что-то не то в качестве записи таблицы перехода.
И пользователь мог подумать, что это дает какой-то элемент в таблице, а на деле его нет.
Потому в этом случае лучше сообщить об ошибке.
Это можно сделать так: составим список всех записей, выкинем из него список из отобранных записей и проверим результат на пустоту.
После определения \verb"s_data" и \verb"e_data" можно добавить следующий код
\begin{cppcode}
  using TFullList = TL::RemoveCVRec<std::remove_cv_t<decltype(data)>>;
  using TSList = TL::RemoveCVRec<std::remove_cv_t<decltype(s_data)>>;
  using TEList = TL::RemoveCVRec<std::remove_cv_t<decltype(e_data)>>;
  using TSEList = TL::Merge2<TSList, TEList>;
  using TUnList = TL::Diff<TFullList, TSEList>;

  static_assert(TL::isEmpty<TUnList>, "Unsupported entry of transition table!");
\end{cppcode}
Сами списки мы получаем как тип соответствующего \verb"Tuple", который мы получаем с помощью \verb"decltype".
Дальше мы рекурсивно убираем все квалификаторы \verb"const/volatile", которые могли объявиться из-за \verb"constexpr" контекста.
И потом находим разность между всеми записями и отобранными.
Тут же можно сделать печать списка, если он не пуст, чтобы сообщить пользователю, какие именно записи были лишними.
Но я не стал так заморачиваться.

\paragraph{Получение списка всех возможных состояний}

Для создания State Machine нужно иметь полный список состояний.
Он спрятан в двух списках, которые мы построили в прошлом разделе.
Теперь надо вытащить все состояния.
Для ребер это их начала и концы, для записей вершин это сами вершины.
Как обычно используем имплементацию для проверки корректности данных
\begin{cppcode}
template<class Data, bool isCorrect = /*checks*/>
struct ExtractStatesImpl {
  static_assert(
      TL::isList<Data>, "The argument of ExtractStates MUST be a list of transition data!");
  ...
};
\end{cppcode}
После чего для корректных данных имплементируем вычленение состояний:
\begin{cppcode}
template<class Data>
struct ExtractStatesImpl<Data, true> {
  using StList = StateList<Data>;
  using EdgeList = EdgeList<Data>;
  
  using StStates = TL::NoDuplicates<TL::Map<DSL::StateOp::State, StList>>;
  using EdgeSources = TL::NoDuplicates<TL::Map<DSL::EdgeOp::Source, EdgeList>>;
  using EdgeTargets =
      TL::EraseAll<TL::NoDuplicates<TL::Map<DSL::EdgeOp::Target, EdgeList>>,
                   void>;
  using type = TL::NoDuplicates<TL::Merge<StStates, EdgeSources, EdgeTargets>>;
};
\end{cppcode}
Код в целом все делает в лоб.
Можно было бы избавление от дублей сделать в самом конце и не делать на промежуточных списках, но это не так принципиально.
И теперь публичная часть:
\begin{cppcode}
template<class TList>
struct ExtractStates_t : detail::ExtractStatesImpl<TList> {};

template<class TList>
using ExtractStates = typename ExtractStates_t<TList>::type;
\end{cppcode}


\paragraph{Получение списка начальных состояний}

Теперь нужно выделить из всех состояний только начальные и только конечные.
Я пропущу проверку корректности и сразу разберем специализацию для проверенных аргументов:
\begin{cppcode}
template<class TList>
struct ExtractInitialStatesImpl<TList, true> {
  using StList = EnExList<TList>;
  using EdgeList = EdgeList<TList>;

  using InitStList = TL::NoDuplicates<TL::Map<
      DSL::StateOp::State,
      TL::Filter<DSL::StateOp::Bind<Core::Initial>::template is_t, StList>>>;

  using InitSourceList = TL::NoDuplicates<
      TL::Map<DSL::EdgeOp::Source,
              TL::Filter<DSL::EdgeOp::Bind<Core::Initial>::template isSource_t,
                         EdgeList>>>;
  using InitTargetList = TL::NoDuplicates<
      TL::Map<DSL::EdgeOp::Target,
              TL::Filter<DSL::EdgeOp::Bind<Core::Initial>::template isTarget_t,
                         EdgeList>>>;

  using type =
      TL::NoDuplicates<TL::Merge<InitStList, InitSourceList, InitTargetList>>;
};
\end{cppcode}
И публичная часть для пользователя:
\begin{cppcode}
template<class TList>
struct ExtractInitialStates_t : detail::ExtractInitialStatesImpl<TList> {};

template<class TList>
using ExtractInitialStates = typename ExtractInitialStates_t<TList>::type;
\end{cppcode}
Вычленение конечных состояний выполняется точно так же.
Потому не буду его тут расписывать.

\paragraph{Получение списка всех возможных событий}

События встречаются только в списке ребер.
Потому нам не нужны записи о вершинах.
И я так же пропущу общую версию шаблона с проверками:
\begin{cppcode}
template<class Edges>
struct ExtractEventsImpl<Edges, true> {
  using type =
      TL::EraseAll<TL::NoDuplicates<TL::Map<DSL::EdgeOp::Event, Edges>>, void>;
};
\end{cppcode}
И публичная часть кода:
\begin{cppcode}
template<class Edges>
struct ExtractEvents_t : detail::ExtractEventsImpl<Edges> {};

template<class Edges>
using ExtractEvents = typename ExtractEvents_t<Edges>::type;
\end{cppcode}

\paragraph{Вырезаем статусы из ребер}

Теперь когда мы знаем статусы всех вершин, потому что знаем полный список вершин и списки начальных и конечных вершин, мы можем вырезать эти статусы из всех ребер.
В начале нужен метод для обработки одного ребра:
\begin{cppcode}
namespace detail {
template<class Edge>
struct TrimEdge_t : TL::Slice_t<Edge, EL::List<int, 0, 2, 3, 4, 5>> {};
template<class Edge>
using TrimEdge = typename TrimEdge_t<Edge>::type;
\end{cppcode}
То есть мы игнорируем вторую и последнюю запись в ребре.
Теперь надо ишь применить эту операцию ко всем ребрам.
Если пропустить проверки, то содержательный код выглядит так:
\begin{cppcode}
template<class Edges>
struct TrimEdgesImpl<Edges, true> {
  using type = TL::Map<TrimEdge, Edges>;
};
\end{cppcode}
И публичная часть для пользователя
\begin{cppcode}
template<class Edges>
struct TrimEdges_t : detail::TrimEdgesImpl<Edges> {};

template<class Edges>
using TrimEdges = typename TrimEdges_t<Edges>::type;
\end{cppcode}

\paragraph{Склеиваем и прореживаем записи для входных/выходных данных}

С записями вершин пока что хаос.
Синтаксис FrontEnd-а позволяет следующие записи
\begin{cppcode}
*S
on_entry >>= S
on_entry >>= S >>= on_exit
S >>= on_exit
\end{cppcode}
В итоге у нас может содержаться аж $4$ записи приходящиеся на одно состояние.
При перепаковке мы теперь имеем $4$ записи вида
\begin{cppcode}
TL::List<Value_t<A,  do_nothing>, S, Initial, Value_t<A,  do_nothing>>
TL::List<Value_t<En, on_entry>,   S, Regular, Value_t<A,  do_nothing>>
TL::List<Value_t<En, on_entry>,   S, Regular, Value_t<Ex, on_exit>>
TL::List<Value_t<A,  do_nothing>, S, Regular, Value_t<Ex, on_exit>>
\end{cppcode}
Нам надо проверить, что среди всех таких записей с одинаковым \verb"S" на входное действие указано не более одного нетривиального действия, и аналогично для выходного действия.
Потому нам надо склеить эти $4$ записи в одну и отрезать статус вершины:
\begin{cppcode}
TL::List<Value_t<En, on_entry>, S, Value_t<Ex, on_exit>>
\end{cppcode}
Поэтому все танцы с бубном ниже будут связаны с попыткой произвести все указанные действия корректно.
Кроме того, я хочу, чтобы в этой таблице для каждого состояния \verb"S" присутствовала в точности одна запись.
А потому для отсутствующих состояний надо будет добавить записи с тривиальными действиями на вход и выход.
Это было не обязательно делать, но это упростит работу с кодом внутри машины состояний.
Я опять проигнорирую проверку условий корректности и перейду к сути имплементации:
\begin{cppcode}
template<class EnExList, class States>
struct ExtractEntryExitImpl<EnExList, States, true> {
  template<class Node>
  using IsOnEntryNonTrivial =
      Logic::Not_w<DSL::StateOp::isOnEntryTrivial_t<Node>>;
  template<class S>
  using OnEntryList = TL::NoDuplicates<
      TL::Map<DSL::StateOp::OnEntryW,
              TL::Filter<IsOnEntryNonTrivial,
                         TL::Filter<DSL::StateOp::Bind<S>::template isSame_t,
                                    EnExList>>>>;
  using Empty = TL::EmptyLike<EnExList>;
  using LNoAction = TL::MakeLike<NoAction, Empty>;

  using OnEntryActions = TL::Map<OnEntryList, States>;

  template<class S>
  using IsOnExitNonTrivial = Logic::Not_w<DSL::StateOp::isOnExitTrivial_t<S>>;
  template<class St>
  using OnExitList = TL::NoDuplicates<
      TL::Map<DSL::StateOp::OnExitW,
              TL::Filter<IsOnExitNonTrivial,
                         TL::Filter<DSL::StateOp::Bind<St>::template isSame_t,
                                    EnExList>>>>;
  using OnExitActions = TL::Map<OnExitList, States>;

  template<class List>
  using AddTrivialAction = IfElse_u<TL::isEmpty<List>, Type_t<NoAction>, TL::Front_t<List>>;
  
  using OnEntryActionList = TL::Map<AddTrivialAction, OnEntryActions>;
  using OnExitActionList = TL::Map<AddTrivialAction, OnExitActions>;
  
  using type = TL::Map<TL::Join, TL::Zip<TL::Zip<OnEntryActionList, States>,
                                         TL::Map<TL::List, OnExitActionList>>>;
};
\end{cppcode}
% TO DO
% Надо разобраться с типами списков в точности!
Давайте разберемся, что тут происходит:
\begin{enumerate}
\item В начале заводим предикат \verb"IsOnEntryNonTrivial", чтобы отобрать все записи, где есть нетривиальное действие на вход.

\item Мы хотим для каждого состояния \verb"S" построить список из записей, где действие \verb"on_entry" не тривиальное.
Так как \verb"S" у нас неизвестное и будет пробегать весь список \verb"States", то мы задаем шаблон \verb"OnEntryList" для такого действия.

\item Давайте разбираться как работает шаблон:
\begin{cppcode}[\nonumbers]
template<class S>
using OnEntryList = TL::NoDuplicates<
    TL::Map<DSL::StateOp::OnEntryW,
            TL::Filter<IsOnEntryNonTrivial,
                       TL::Filter<DSL::StateOp::Bind<S>::template isSame_t, EnExList>>>>;
\end{cppcode}
Внутри мы делаем \verb"Bind<S>" и вызываем оператор сравнения с \verb"S".
Операция \verb"Filter" таким образом отбирает все вершины, с состоянием \verb"S".
После чего следующая операция \verb"Filter" отбирает среди таких записей те, что содержат нетривиальное \verb"on_entry" действие.
После чего операция \verb"Map" применяет ко всем полученным записям \verb"OnEntryW" вычленяя тип вида \verb"Value_t<En, on_entry>".
В конце мы избавляемся от дубликатов.
И в этом месте стоит проверить, что размер полученного списка не более чем $1$, но я пожалуй пропущу это.

\item Далее мы просто применяем \verb"OnEntryList" для всех возможных состояний.
В итоге у нас для каждого состояния \verb"S" получается список из записей вида
\verb"Value_t<En, on_entry>".

\item Дальше идет аналогичное выделение списка нетривиальных действий на выход, которые складываются в \verb"OnExitActions".

\item Теперь обсудим, что мы сейчас имеем.
\begin{cppcode}[\nonumbers]
OnEntryActions = TL::List<OnEntryActionS1, OnEntryActionS2, ...>
OnExitActions =  TL::List<OnExitActionS1, OnExitActionS2, ...>
// where
OnEntryActionS1 = TL::List<Value_t<En1, on_entry1>,
                           Value_t<En2, on_entry2>,
                           ... >
...
OnExitActionS1 =  TL::List<Value_t<Ex1, on_exit1>,
                           Value_t<Ex2, on_exit2>,
                           ... >
...
\end{cppcode}
Здесь внешний список \verb"OnEntryActions" имеет такой же размер как \verb"States" и для $k$-ого состояния \verb"Sk" он хранит список всех заданных пользователем действий вида \verb"on_entry".
Внутренний список всех действий для вершины \verb"Sk" я обозначил \verb"OnEntryActionSk".
И для каждого такого списка может быть три случая:
\begin{cppcode}[\nonumbers]
1) OnEntryActionS = TL::List<>
2) OnEntryActionS = TL::List<Value_t<En, on_entry>>
3) OnEntryActionS = TL::List<Value_t<En1, on_entry1>,
                             Value_t<En2, on_entry2>,
                             ... >
\end{cppcode}
Первый случай означает, что пользователь не задал никакого нетривиального действия.
Второй означает, что пользователь задал одно единственное нетривиальное действие.
И последний случай пользователь задал более одного нетривиального действия.
В последнем случае надо сообщать об ошибке.%
\footnote{Но я позволю себе проигнорировать этот замечательный факт.}
А в первом надо добавить тривиальное действие.
Для чего мы заводим вспомогательный шаблон \verb"AddTrivialAction".
\begin{cppcode}[\nonumbers]
template<class List>
using AddTrivialAction = IfElse_u<TL::isEmpty<List>, Type_t<NoAction>, TL::Front_t<List>>;
\end{cppcode}

\begin{cppcode}[\nonumbers]
using OnEntryActionList = TL::Join<TL::Map<AddTrivialAction, OnEntryActions>>;
using OnExitActionList = TL::Join<TL::Map<AddTrivialAction, OnExitActions>>;
\end{cppcode}
Здесь нужен внешний \verb"Join" чтобы избавиться от списка списков.
То есть мы переводим
\begin{cppcode}[\nonumbers]
OnEntryActions = TL::List<TL::List<Value_t<En1, on_entry1>>,
                          TL::List<Value_t<En2, on_entry2>>,
                          ... >
\end{cppcode}
в
\begin{cppcode}[\nonumbers]
OnEntryActionList = TL::List<Value_t<En1, on_entry1>,
                             Value_t<En2, on_entry2>,
                             ... >
\end{cppcode}
При этом тут количество записей такое же как в \verb"States" и каждая позиция соответствует правильному действию для данного \verb"State".
Аналогично для \verb"OnExitActionList".

\item Теперь слияние.
\begin{cppcode}[\nonumbers]
using type = TL::Map<TL::Join, TL::Zip<TL::Zip<OnEntryActionList, States>,
                                       TL::Map<TL::List, OnExitActionList>>>;
\end{cppcode}
Давайте я поясню на примере, что тут конкретно происходит.
Пусть у нас всего $2$ вершины и списки выглядят так:
\begin{cppcode}[\nonumbers]
States            = TL::List<S0, S1>
OnEntryActionList = TL::List<Values_t<A0, on_entry0>, Value_t<A1, on_entry1>>
OnExitActionList  = TL::List<Values_t<A0, on_exit0>,  Value_t<A1, on_exit1>>
\end{cppcode}
Причем длины списков у нас одинаковые и теперь мы хотим слить покомпонентно все списки.
\begin{cppcode}[\nonumbers]
TL::List<
    TL::List<Values_t<A0, on_entry0>, S0, Values_t<A0, on_exit0>>,
    TL::List<Values_t<A1, on_entry1>, S1, Values_t<A1, on_exit1>>
>
\end{cppcode}
По хорошему мы хотим \verb"Zip" для трех списков.
Я поленился его писать, а потому будем писать руками.
Последовательность операций можно визуализировать так:
\begin{cppcode}[\nonumbers]
En [f0, f1] |   
St [S0, S1] |=> [[f0, S0], [f1, S1]] |=> [[[f0, S0], [g0]], [[f1, S1], [g1]]] |=>
Ex [g0, g1] |   [[g0], [g1]]         |

=> [[f0, S0, g0], [f1, S1, g1]]
\end{cppcode}
В начале мы к паре \verb"[En, St]" применяем \verb"Zip", а на элементы \verb"Ex" навешиваем список.
Потом делаем еще раз \verb"Zip" и получаем список пар.
И теперь надо в каждой паре внутри стереть внутренние скобки, что делает поэлементное применение \verb"Join".
\end{enumerate}
И теперь пользовательская часть кода:
\begin{cppcode}
template<class TList, class States>
struct ExtractEntryExit_t : detail::ExtractEntryExitImpl<TList, States> {};

template<class TList, class States>
using ExtractEntryExit = typename ExtractEntryExit_t<TList, States>::type;
\end{cppcode}

\subsection{BackEnd Генерация машины по таблице}


Теперь мы подготовили все вспомогательные методы для конвертации данных в списки типов и можно строить машину состояний.
Я начну сверху вниз.
В начале покажу как выглядит ее интерфейс и методы снаружи, и уже потом буду добавлять имплементацию.

\subsubsection{Интерфейс State Machine}

Глобально код создающий машину состояний будет выглядеть так.
Внутри служебного пространства имен заведем класс для имплементации.
Как обычно разделенный на две части: общая версия шаблона для сигнализации об ошибках в данных и специализация для корректных данных:
\begin{cppcode}
namespace Implementation {
template<class Edges, class TEnExList, class InitStates, class TFinalStates,
         class TEvents, bool isCorrect = true>
class SM {
  /*error messages*/
};

template<class Edges, class TEnExList, class InitStates, class TFinalStates,
         class TEvents>
class SM<Edges, TEnExList, InitStates, TFinalStates, TEvents, true> {
  /*Implementation*/
};
} // namespace Implementation
\end{cppcode}
Теперь нам нужен вспомогательный шаблон для конвертации пользовательской таблицы с описанием State Machine в данные, по которым строится имплементация:
\begin{cppcode}
namespace Implementation {
template<class SmTable>
struct SmFromTable_t {
  using Data = Convert::ExtractTransitionData<SmTable>;

  using EdgeList = Convert::EdgeList<Data>;
  using EnExList = Convert::EnExList<Data>;
  using States = Convert::ExtractStates<Data>;
  using InitStates = Convert::ExtractInitialStates<Data>;
  using FinalStates = Convert::ExtractFinalStates<Data>;
  using Events = Convert::ExtractEvents<EdgeList>;
  using Edges = Convert::TrimEdges<EdgeList>;
  using EnExListT = Convert::ExtractEntryExit<EnExList, States>;

  using type = SM<Edges, EnExListT, InitStates, FinalStates, Events>;
};

template<class SmTable>
using SmFromTable = typename SmFromTable_t<SmTable>::type;
} // namespace Implementation
\end{cppcode}
Здесь мы по-порядку применяем те методы, которые обсуждались в предыдущем разделе.
\begin{cppcode}
template<class SmTable>
class StateMachine : public Implementation::SmFromTable<SmTable> {
  using Base = Implementation::SmFromTable<SmTable>;
  using Base::Base;
};
\end{cppcode}
Публичное наследование от имплементации.
При этом открываются все публичные методы, однако, конструкторы таким образом не наследуются.
А потому тут я дополнительно пробрасываю конструкторы.
Теперь все что нужно -- обсудить как внутри устроена имплементация для шаблона \verb"SM".

\subsubsection{Имплементация по модулю визитеров}

\paragraph{Интерфейс имплементации}

Вот как выглядит интерфейс имплементации:
\begin{cppcode}
template<class TEdges, class TEnExList, class TStates, class InitStates, class TFinalStates,
         class TEvents, bool isCorrect = /*checks*/>
class SM {
  /*error messages*/
};

template<class TEdges, class TEnExList, class TStates, class InitStates, class TFinalStates,
         class TEvents>
class SM<TEdges, TEnExList, TStates, InitStates, TFinalStates, TEvents, true> {
public:
  /* using declarations */

  /* ctors, assignments and dtors */

  template<class Event>
  bool process(Event&& event);

  template<class State>
  bool isAt() const;

  void reset();
  bool start();

  bool isComplete() const ;
  bool isStarted() const;
  bool isRunning() const;

  template<class Event>
  static constexpr bool acceptEvent() noexcept;
  template<class State>
  static constexpr bool isValidState() noexcept;
  template<class State>
  static constexpr bool isFinalState() noexcept;
  template<class State>
  static constexpr bool isInitialState() noexcept;
  
private:
  /*private methods*/

  AState state_ = Idle{};
};
\end{cppcode}

\paragraph{Определение псевдонимов}

Начать я хочу с определения псевдонимов в самом начале класса.
Это нужно, чтобы подготовить всю статическую информацию для использования.

В начале мы хотим изменить множество вершин \verb"TStates", потому что мы хотим добавить вспомогательные служебные состояния.
\begin{cppcode}[\nonumbers]
using Idle = Back::IdleState; // some empty struct
using Initial = I<TL::Front<InitStates>>; // I<...> is a template defining an empty struct
using FinalStates = TL::Map<F, TFinalStates>; // F<...> is a template defining an empty struct

using States = TStates;
using StatesExt = TL::PushFront<Initial, TL::Merge2<FinalStates, States>>;
\end{cppcode}
Мы добавили состояние \verb"Idle" для незапущенной машины, завели новое служебное начальное состояние и для каждого финального состояния завели новые финальные состояния.
Я оставляю список исходных состояний для пользователя снаружи, а расширенный список \verb"StatesExt" будет использоваться внутри машины.
Теперь нужно добавить ребра для служебных вершин.
Нужно добавить одно ребро из служебной начальной вершины в настоящую начальную вершину.
\begin{cppcode}[\nonumbers]
using InitialEdge = TL::ConvertLike<
    TL::List<Initial, void, NoGuard, NoAction, TL::Front<InitStates>>, TEdges>;
\end{cppcode}
Метод \verb"ConvertLike" нужен для того, чтобы конвертировать построенный список из \verb"TL::List" в тот вид списка, который используется в \verb"TEdges".
И для каждой финальной вершины сделать ребро в соответствующую служебную вершину.
\begin{cppcode}[\nonumbers]
template<class S>
using FinalEdge = TL::ConvertLike<TL::List<F<S>, void, NoGuard, NoAction, S>, Edges>;

using FinalEdges = TL::Map<FinalEdge, TFinalStates>;
\end{cppcode}
И теперь все объединяем вместе:
\begin{cppcode}[\nonumbers]
using Edges = TL::PushFront<InitialEdge, TL::Merge2<FinalEdges, TEdges>>;
\end{cppcode}
Теперь надо добавить соответствующие действия на вход и выход для служебных вершин
\begin{cppcode}[\nonumbers]
template<class S>
using FinalEnEx = TL::ConvertLike<TL::List<NoAction, S, NoAction>, TEdges>;

using InitEnEx = FinalEnEx<Initial>;
using FinalEnExList = TL::Map<FinalEnEx, FinalStates>;
\end{cppcode}
И теперь добавляем все в \verb"TEnExList" так:
\begin{cppcode}[\nonumbers]
using EnExListExt = TL::PushFront<InitEnEx, TL::Merge2<FinalEnExList, TEnExList>>;
\end{cppcode}
Теперь определим тип \verb"AState" для хранения произвольного состояния.
У нас это будет \verb"std::variant" от всех вершин в \verb"StatesExt".
Это делается так:
\begin{cppcode}[\nonumbers]
using AState = TL::Convert<StatesExt, std::variant>;
\end{cppcode}
И под конец заводим псевдонимы для событий.
Так как я не завожу никаких служебных событий, то доступные события определяются просто:
\begin{cppcode}[\nonumbers]
using Events = TEvents;
\end{cppcode}
Важно отметить, что мы храним все типы состояний и событий нормализованными, то есть к ним предварительно была применена \verb"std::remove_cvref_t".%
\footnote{В имплементации у нас вызывалось рекурентное избавление от \texttt{const/volatile} квалификаторов, и одна из проверок на корректность должна быть в том, чтобы элементы \texttt{States} и \texttt{Events} обязаны быть не ссылками.}

\paragraph{Имплементация информационных методов}

Начнем с имплементации статических служебных методов.
\begin{cppcode}
template<class Event>
static constexpr bool acceptEvent() noexcept {
  using PEvent = std::remove_cvref_t<Event>;
  return TL::Contain<Events, PEvent>;
}

template<class State>
static constexpr bool isValidState() noexcept {
  using PState = std::remove_cvref_t<State>;
  return TL::Contain<States, PState>;
}

template<class State>
static constexpr bool isFinalState() noexcept {
  using PState = std::remove_cvref_t<State>;
  return TL::Contain<FinalStates, PState>;
}

template<class State>
static constexpr bool isInitialState() noexcept {
  using PState = std::remove_cvref_t<State>;
  return areSame2<Initial, PState>;
}
\end{cppcode}
Теперь обсудим простые методы сообщающие информацию о статусе машины:
\begin{cppcode}
bool isComplete() const {
  return std::visit(Visitors::IsComplete<FinalStates>{}, state_);
}

bool isStarted() const {
  return !std::holds_alternative<Idle>(state_);
}

bool isRunning() const {
  return isStarted() && !isComplete();
}
\end{cppcode}
Обратите внимание, что для метода \verb"isStarted" мы используем STL функцию.
Ее важное условие -- тип \verb"Idle" обязан быть среди типов указанных в \verb"std::variant".
Но у нас это условие выполняется.
А вот для выяснения находимся ли мы в одном из финальных состояний уже придется писать своего визитора.
Этот объект будет вынимать текущее состояние из \verb"std::variant" и уже зная что там лежит выполнять нужное действие.
Я пока отложу вопрос имплеменатции визитеров.
На них будет рассчитан следующий раздел, а здесь я их использую как черные ящики.
Нужно лишь обратить внимание какую информацию визитер от меня получает.
И мы видим, что
\begin{cppcode}
Visitors::IsComplete<FinalStates>{}
\end{cppcode}
получает только список финальных состояний и никакой run-time информации.
Еще один метод проверки состояния:
\begin{cppcode}
template<class State>
bool isAt() const {
  using PState = std::remove_cvref_t<State>;
  return std::visit(Visitors::IsAt<PState>{}, state_);
}
\end{cppcode}
Тут просто передаем работу соответствующему визитеру.

\paragraph{Запуск и остановка машины}

Тут я покажу, как выглядят методы, которые переключают машину из \verb"Idle" состояния в начальное и из любого в \verb"Idle".
\begin{cppcode}
bool start() {
  if constexpr (isValidState<Idle>()) {
    forceMoveTo<Idle>();
    processAnonymous<Idle>();
    return true;
  } else {
    return false;
  }
}
\end{cppcode}
И вот тут мы видим два приватных вспомогательных метода.
Первый перемещает машину в состояние \verb"Idle", а второй выполняет все анонимные переходы начиная из текущего состояния.
В целом можно было бы не указывать \verb"Idle" во втором методе, но так я просто экономлю один \verb"switch", ведь я знаю, кто сейчас лежит внутри \verb"std::variant".

Для остановки машины используем:
\begin{cppcode}
  void reset() {
    forceMoveTo<Idle>();
  }
\end{cppcode}
Предполагается, что перемещение машины из вершину в вершину насильно не вызывает действий \verb"on_exit" и \verb"on_entry".
Однако, часто при выходе из вершины \verb"on_exit" играет роль деструктора, который чистит данные.
А потому \verb"forceMoveTo" при выходе из текущей вершины будет вызывать \verb"on_exit" для данной вершины.
В целом операция сброса состояния все равно может оставить данные в некорректном состоянии.
Но что точно верно: такой метод можно использовать в служебных состояниях.
Например сброс всегда корректно сделать из финального состояния в \verb"Idle".

\paragraph{Имплементация обработчика событий}

Давайте запишем, как будет выглядеть обработчик событий:
\begin{cppcode}
template<class Event>
bool process(Event&& event) {
  if (!isRunning())
    return false;
  using PEvent = std::remove_cvref_t<Event>;
  if constexpr (acceptEvent<PEvent>()) {
    if (!processOnce(std::forward<Event>(event))) {
      return false;
    }
    processAnonymous();
    return true;
  } else {
    return false;
  }
}
\end{cppcode}
Тут происходит проверка на то, что машина в рабочем состоянии, иначе любое сообщение игнорируется.
После мы чистим тип сообщения от ссылок и \verb"const/volatile" модификаторов и смотрим, принимается ли это сообщение машиной.
Если не принимается, то просто возвращаем \verb"false".
А если принимается, то обработка идет по следующему правилу.
\begin{enumerate}
\item В начале мы смотрим можно ли пройти хотя бы по одному ребру с указанным типом сообщения.
Если нельзя, то мы возвращаем \verb"false" и выходим.

\item Если можно, то мы оказываемся в новом состоянии (вообще говоря) и из него мы теперь выполняем все возможные анонимные переходы.
\end{enumerate}

\paragraph{Приватные методы и их имплементации}

Давайте я в начале напишу все приватные методы, которые у меня есть:
\begin{cppcode}
template<class Event>
bool processOnce(Event&& event);

template<class State>
bool forceMoveTo();

template<class State>
bool processAnonymousOnce();
bool processAnonymousOnce();

template<class State>
bool processAnonymous();
bool processAnonymous();
\end{cppcode}
Все методы возвращают \verb"true", если они смогли совершить переход и \verb"false" иначе.
Обратите внимание, что для анонимных переходов есть две версии: шаблонная и не шаблонная.
Шаблонная нужна для случаев, когда мы уже знаем в каком состоянии какого типа находится машина, чтобы избежать лишний \verb"switch" по текущему состоянию.
Напишем имплементации по порядку.
Следующий метод делает один шаг вдоль ребра перехода, если найдется подходящее ребро.
\begin{cppcode}
template<class Event>
bool processOnce(Event&& event) {
  assert(isRunning());
  using PEvent = std::remove_cvref_t<Event>;
  return std::visit(
             Visitors::ProcessOnce<Edges, EnExListExt, AState, PEvent>(state_, event), state_);
}
\end{cppcode}
Тут я так же свожу работу к соответствующему визитеру.
Обратите внимание, что он запоминает текущее состояние \verb"state_" в нераспакованном виде и сообщение.
Состояние в нераспакованном виде нужно для того, чтобы его можно было поменять.
Так же мы передаем в виде шаблонной информации все необходимые для перехода данные.
Следующий метод для насильного перемещения машины.
Данный метод делает только перемещение, но не выполняет последующих операций, если таковые имеются.
\begin{cppcode}
template<class State>
bool forceMoveTo() {
  static_assert(areSame2<State, std::remove_cvref_t<State>>, "Internal error!");
  std::visit(Visitors::ForceMove<Edges, EnExListExt, AState, State>(state_), state_);
  return true;
}
\end{cppcode}
Теперь метод для выполнения одного анонимного перехода.
\begin{cppcode}
template<class State>
bool processAnonymousOnce() {
  static_assert(areSame2<State, std::remove_cvref_t<State>>, "Internal error!");
  assert(std::get_if<State>(&state_) && "Internal error!");
    return Visitors::AnonymousTransitOnce<Edges, EnExListExt, AState>(state_)
        .template call<State>();
}

bool processAnonymousOnce() {
    return std::visit(Visitors::AnonymousTransitOnce<Edges, EnExListExt, AState>(state_), state_);
}
\end{cppcode}
И обработчик всех анонимных сообщений:
\begin{cppcode}
template<class State>
bool processAnonymous() {
  static_assert(areSame2<State, std::remove_cvref_t<State>>, "Internal error!");
  if (!processAnonymousOnce<State>())
    return false;
  while (processAnonymousOnce()) {
  }
  return true;
}

bool processAnonymous() {
  if (!processAnonymousOnce())
    return false;
  while (processAnonymousOnce()) {
  }
  return true;
}
\end{cppcode}
Такой обработчик пробует сделать один шаг (шаблонная версия при этом делает его из известного состояния).
И если получилось, то потом делаем в цикле шаги до тех пор, пока хотя бы один шаг был сделан.

\subsubsection{Обеспечение семантики}

На мой взгляд это настолько важный вопрос, что для него полагается отдельный раздел.

\paragraph{Создание машины}

Классически мы считаем, что машина создается в незапущенном состоянии \verb"Idle".
Потому дефолтный конструктор для машины состояний можно использовать по умолчанию:
\begin{cppcode}
template<class TEdges, class TEnExList, class TStates, class TInitStates,
         class TFinalStates, class TEvents>
class SM<TEdges, TEnExList, TStates, TInitStates, TFinalStates, TEvents, true> {
public:
  SM() = default;
  ...
private:
  AState state_ = Idle{};
};
\end{cppcode}
Если при этом машина поддерживает данные внутри себя, то их так же можно сообщить в конструкторе
\begin{cppcode}
Data data;
StateMachine sm(std::move(data), 1, 1.2);
\end{cppcode}
Тип этих данных можно вытащить из аргументов \verb"Guard"-ов и \verb"Action"-ов отличных от служебных данных вроде состояний и событий.
И можно сделать конструктор, который зависит только от этих аргументов и поглощает внутрь себя их начальные состояния.
Вопрос в том, а как распаковать список аргументов.
Давайте поясню, что это значит.
Вот у вас есть класс
\begin{cppcode}
template<class TEdges, class TEnExList, class TStates, class TInitStates,
         class TFinalStates, class TEvents>
class SM<TEdges, TEnExList, TStates, TFinalStates, TEvents, true> {
public:
  using DataList = /*extract from the transition table*/;
  ...
private:
  AState state_ = Idle{};
};
\end{cppcode}
И теперь мы знаем, что \verb"DataList" -- это список вида
\begin{cppcode}[\nonumbers]
DataList = TL::List<Data, int, double>;
\end{cppcode}
И мы хотим создать конструктор вида
\begin{cppcode}[\nonumbers]
SM(Data&&, int&&, double&&);
\end{cppcode}
Как распаковать список, когда мы знаем только его имя?
Один из способов -- добавить список типов \verb"DataList" в число аргументов по умолчанию и специализировать его на распакованный вариант.
\begin{cppcode}
template<class TEdges, class TEnExList, class TStates, class TInitStates,
         class TFinalStates, class TEvents, class DataList = ExtractData<TEdges, TEnExList>,
         bool isCorrect = /*checkes*/>
class SM { /*error messages*/ };

template<class TEdges, class TEnExList, class TStates, class TInitStates,
         class TFinalStates, class TEvents, template<class...> class List, class... Ts>
class SM<TEdges, TEnExList, TStates, TFinalStates, TEvents, List<Ts...>, true> {
public:
  using DataList = List<Ts...>;
  SM(Ts&&... args) : data_(MakeTuple(std::move(args)...)) {}
  ...
private:
  AState state_ = Idle{};
  Tuple<std::remove_cvref_t<Ts>...> data_;
};
\end{cppcode}
В данном случае \verb"Ts&&" является rvalue ссылкой, а не универсальной ссылкой, потому что \verb"Ts" не является шаблонным аргументом для самого конструктора.
На уровне конструктора это уже конкретный известный тип.
При такой имплементации сложно обеспечить различные виды ссылок для каждого аргумента.

Другой вариант -- сделать шаблонный конструктор и надеяться, что данные преобразуются как надо.
\begin{cppcode}
template<class TEdges, class TEnExList, class TStates, class TInitStates,
         class TFinalStates, class TEvents, bool isCorrect = /*checkes*/>
class SM { /*error messages*/ };

template<class TEdges, class TEnExList, class TStates, class TInitStates,
         class TFinalStates, class TEvents>
class SM<TEdges, TEnExList, TStates, TFinalStates, TEvents, true> {
public:
  using DataList = /*extract from the transition table*/;
  using Data = /*create Tuple from DataList*/;
  template<class... Ts, class = std::enable_if_t<(sizeof...(Ts) == TL::Size<DataList>)>>
  SM(Ts&&... args) : data_(std::forward<Ts>(args)...) {}
  ...
private:
  AState state_ = Idle{};
  Data data_;
};
\end{cppcode}
Аналогичным приемом мы пользовались в случае имплементации \verb"Tuple" (см.~\ref{}).

Так же можно делегировать вопрос хранения и создания зависимостей отдельному классу и потом использовать Polici Based Design (или Mixing-и).
Про это я немного писал тут~\ref{}.
Давайте кратко скажу в чем суть.
Давайте посвятим хранению данных отельный объект:
\begin{cppcode}
template<class List>
struct Storage;

template<template<class...> class List, class... Ts>
class Storage<List<Ts...>> {
public:
  Storage(Ts&&... args) : data_(std::move(args)...) {}
  template<class... Args, std::enable_if_t<(sizeof...(Args) == sizeof...(Ts))>>
  Storage(Args&&... args) : data_(std::forward<Args>(args)...) {}
protected:
  Tuple<std::remove_cvref_t<Ts>...> data_;
};
\end{cppcode}
Тут можно выбрать какой из вариантов конструкторов вы хотите использовать: поглощающий или использующий perfect forwarding (см.~\ref{}).
Теперь надо лишь унаследоваться от этого класса:
\begin{cppcode}
template<class TEdges, class TEnExList, class TStates, class TInitStates,
         class TFinalStates, class TEvents>
class SM<TEdges, TEnExList, TStates, TFinalStates, TEvents, true>
    : Storage<ExtractDataList<TEdges, TEnExList>>  {
  using Base = Storage<ExtractDataList<TEdges, TEnExList>>;
public:
  using Base::Base; // inherit constructors
  ...
};
\end{cppcode}
И теперь надо лишь прокинуть конструкторы с помощью \verb"using" директивы.
Так мы отделяем вопрос хранения данных от использования данных.
В частности можно использовать разные стратегии хранения.
Например можно завернуть внутренний \verb"Tuple" в \verb"unique_ptr" или \verb"HeapHolder" (см.~\ref{}) и получить почти бесплатное перемещение данных.
Либо использовать \verb"Tuple" как и написано выше, что дает хорошую локальность данных.

\paragraph{Разрушение машины}

Теперь надо понять, что делать для разрушения машины.
Я не имплементировал поддержку данных в примере для State Machine, однако, машина состояний скорее всего поддерживает данные, которые то инициализируются, то освобождаются, а возможно даже удаляются при переходе между состояниями.
В частности при входе в какое-то состояние у вам может произойти аллокация ресурса, который надо освободить при выходе.
Это означает, что находясь в произвольной вершине, вы  вообще говоря не можете просто так воспользоваться деструктором по умолчанию.
Тем не менее, если машина находится в \verb"Idle" состоянии или в одном из финальных состояний, то дефолтный деструктор должен сработать.
Если же машина находится не в дефолтном состоянии, то нужно выполнить действия по выходу из вершины (что должно освбодить ресурсы захваченные при входе в вершину) и это все.
Так что деструктор будет иметь вид:
\begin{cppcode}
~SM() noexcept {
  reset();
}
\end{cppcode}
где функция \verb"reset" в нашем случае подразумевает внутри \verb"forceMoveTo<Idle>", который по нашему допущению в обход принятого поведения вызывает \verb"on_exit" действие из текущей вершины.

\paragraph{Перемещение}

При перемещении машины
\begin{cppcode}
StateMachine sm0;
StateMachine sm1 = std::move(sm0);
StateMachine sm2;
sm2 = std::move(sm1);
\end{cppcode}
нужно заботиться о том, чтобы оставить предыдущую машину в <<пустом>> состоянии.
Машина содержит по сути два вида данных:
\begin{enumerate}
\item текущее состояние \verb"state_".

\item текущие данные \verb"data_".
\end{enumerate}
И мы должны обеспокоиться о перемещении их всех.
Нам нужно понять состояние машины, в котором она остается после перемещения.
Обратите внимание, что \verb"std::variant" после перемещения содержит объект того же типа, но из которого мувнули данные.
Даже если мы забудем про данные \verb"data_" и их корректность, то подобная процедура может поставить State Machine в некорректное состояние, потому что состояние могло содержать специальным образом синхронизированные данные (подчиненные некому инварианту), чего больше нет.
На мой взгляд правильным решением будет перевести машину в исходное дефолтное состояние после работы дефолтного конструктора.
Это значит, что надо выставить \verb"state_" в \verb"Idle{}".
Обратите внимание, что тут я подразумеваю, что классы состояний обладают value semantics.
Иначе могут быть всякие проблемы.

Для данных \verb"data_" ситуация следующая.
Если они уже обладают value semantics (см.~\ref{}), то данные корректно переместятся в новую машину.
Однако, то что осталось в \verb"data_" после мува из них может не соответствовать состоянию, в котором машина ожидает данные в начале своей работы.
Потому может понадобиться сброс данных до некоторого дефолтного состояния.
Чтобы не было проблем проще всего дефолтное состояние считать состоянием после дефолтного конструктора, а уже содержательную инициализацию выполняет \verb"on_entry" действие.
В любом случае если машина находится в некотором состоянии, то при входе в нее \verb"on_entry" мог захватить ресурсы и положить их в переменную для State Machine.
Потому нам надо выполнить \verb"on_exit" действие для текущего состояния.
Иначе могут быть проблемы.
\begin{cppcode}
SM(SM&& other) noexcept 
    : state_(std::exchange(other.state_, Idle{})),
      data_(std::exchange(other.data_, Default{})) {
}
SM& operator=(SM&& other) noexcept {
  SM tmp = std::move(other);
  std::swap(state_, tmp.state_);
  std::swap(data_, tmp.data_);
  return *this;
}
\end{cppcode}
Здесь \verb"Default" подразумевает инициализацию каким-то дефолтным значением данных.
Если такая реинициализация ненужна, то можно было бы использовать  \verb"std::move" вместо \verb"std::exchange", но только будьте уверены, что это действительно корректно для ваших данных.

% Поместить этот текст в раздел про value semantics
Еще один общий подход заключается в том, чтобы считать, что  машиной просто нельзя пользоваться после мува из нее.
То есть машина как не инициализированная переменная находится в корректном но мусорном состоянии.
Перед использованием такой переменной в ей надо вначале присвоить некоторое значение и уже потом пользоваться.
При таком подходе надо быть аккуратным, потому что компилятор вам почти никогда не подскажет, что вы мувнули из объекта и им нельзя пользоваться.
Самая главная проблема -- деструктор.
Например если вы мувнули данные из \verb"std::vector"-а выполнив просто \verb"std::memcpy" для локальных данных, то старый вектор содержит все те же самые указатели на те же самые данные.
Использование старого вектора конечно должно считаться UB, однако, как деструктор старого вектора должен понять, что ему не надо исполняться?
Обычно он понимает это потому, что указатели равны \verb"nullptr", а в этом случае только компилятор может уследить за подобным.
Но в случае нескольких единиц трансляции, компилятор просто может не знать, что из объекта мувнули в коде из другой единицы трансляции.
Это все делает подобные подходы в C++ очень опасными и сложно реализуемыми.

\paragraph{Можно ли копировать?}

Если и классы состояний и классы для данных подчиняются value semantics, то ничто не мешает использовать дефолтные конструктор и оператор копирования, потому что и состояние машины и данные будут в корректном и одинаковом состоянии после копирования и не будут влиять друг на друга никаким образом.
\begin{cppcode}
SM(const SM&) = default;
SM& operator=(const SM& other) = default;
\end{cppcode}

\subsubsection{Имплементация визитеров}

Теперь перейдем к имплементации визитеров.
Я для удобства спрячу их в общее пространство имен \verb"Visitors".
Давайте я в начале перечислю всех визитеров, что у нас есть:
\begin{cppcode}
namespace Visitors {
template<class State>
struct IsAt;

template<class FinalStates>
struct IsComplete;

template<class Edges, class StatesEnEx, class AState, class From>
struct AnonymousTransitOnce;

template<class Edges, class StatesEnEx, class AState, class Target>
struct ForceMove;

template<class Edges, class StatesEnEx, class AState, class Event>
struct ProcessOnce;
} // namespace Visitors 
\end{cppcode}
Сразу скажу, что я буду игнорировать данные в State Machine при имплементации визитеров.
Это не принципиально, они легко добавляются в код ниже.
Возможно я когда-нибудь допишу этот код.

\paragraph{Визитер \texttt{IsAt}}

Самый простой визитер для проверки состояния машины:
\begin{cppcode}
template<class State>
struct IsAt {
  static_assert(areSame2<State, std::remove_cvref_t<State>>);
  template<class CState>
  bool operator()(CState&) const {
    using PState = std::remove_cvref_t<CState>;
    return areSame2<State, PState>;
  }
};
\end{cppcode}
Прежде всего я предполагаю, что шаблонные аргументы передаются во все визитеры в нормализованном виде, то есть после применения \verb"std::remove_cvref_t".
Однако, аргумент оператора может не удовлетворять этому условию, потому что \verb"state_" мог быть константным, так как вызвали на константном объекте.
И потому \verb"CState" мне надо сначала почистить от возможного наличия ссылки или \verb"const/volatile" квалификаторов.

\paragraph{Визитер \texttt{IsComplete}}

Тут имплементация почти такая же простая, как и для шаблона \verb"IsAt":
\begin{cppcode}
template<class FinalStates>
struct IsComplete {
  template<class CState>
  bool operator()(CState&) const {
    using PState = std::remove_cvref_t<CState>;
    return TL::Contains<FinalStates, PState>;
  }
};
\end{cppcode}
Так как финальных состояний много, то мы проверяем принадлежность списку \verb"FinalStates".
Тут чистка никакая не нужна, потому что я не ожидаю тут ссылку или указатель, я ожидаю именно список типов.
А все операции для списка типов их вначале чистят от \verb"const/volatile" квалификаторов, так что все и так должно работать, либо надо сообщать об ошибке использования.

\paragraph{Визитер \texttt{AnonymousTransitOnce}}

Вот как выглядит имплементация:
\begin{cppcode}
template<class Edges, class EnExList, class AState>
struct AnonymousTransitOnce {
  AnonymousTransitOnce(AState& state) : state_(state) {
    // need to capture sm.data
  }

  template<class Source>
  bool operator()(Source& state) const {
    using PSource = std::remove_cvref_t<Source>;
    return call<PSource>();
  }

  template<class From>
  bool call() const {
    static_assert(areSame2<From, std::remove_cvref_t<From>>);
    return Back::TransitionFrom<AState, From, void, Edges, EnExList>(state_);
  }
  AState& state_;
};
\end{cppcode}
Визитер перехватывает текущее состояние State Machine по неконстантной ссылке.
Если в машине есть данные, то надо перехватить и их.
Теперь у данного визитера кроме стандартного оператора \verb"operator()" есть еще служебный метод \verb"call".
Суть в том, что \verb"operator()" нужен для распаковки \verb"std::visitor", а \verb"call" знает напрямую какой тип элемента внутри \verb"std::visitor".
Этот метод нужен чтобы избежать лишнего \verb"switch" по вариантам, если снаружи мы и так знаем в каком мы сейчас состоянии (например после \verb"forceMoveTo").

В этой имплементации оператор \verb"operator()" зовет метод \verb"call" после определения типа состояния.
А вот уже оператор \verb"call" вызывает метод
\begin{cppcode}[\nonumbers]
Back::TransitionFrom<AState, From, void, Edges, EnExList>
\end{cppcode}
из BackEnd-а.
Именно этот метод выполняет переход по одному анонимному ребру, если такой переход возможен.
Шаблоны выше является указателем на функцию.
Сейчас для нас эта функция является черным ящиком, но давайте поясним смысл третьего аргумента \verb"void".
Данный аргумент должен указать тип сообщения \verb"Event".
Но для анонимных переходов мы будем использовать \verb"void" для типа сообщения.
Таким образом функция
\begin{cppcode}[\nonumbers]
Back::TransitionFrom<AState, From, Event, Edges, EnExList>
\end{cppcode}
умеет делать переход из состояния \verb"From" по сообщению типа \verb"Event", зная список переходов \verb"Edges" и списк действий на вход и выход для состояний \verb"EnExList".
А в данном визитере мы просто используем частный случай этой функции.
Обратите внимание, что визитер у меня не занимается содержательной работой, он лишь занимается распаковкой \verb"std::visitor".
А содержательную работу как раз и выполняет шаблон \verb"TransitionFrom", которому визитер делегирует эту работу.

\paragraph{Визитер \texttt{ForceMove}}

Имплементация визитера:
\begin{cppcode}
template<class Edges, class EnExList, class AState, class Target>
struct ForceMove {
  static_assert(areSame2<Target, std::remove_cvref_t<Target>>);
  ForceMove(AState& state) : state_(state) {
  }

  template<class Source>
  bool operator()(Source&) const {
    using PSource = std::remove_cvref_t<Source>;
    return Back::ForceTransition<AState, PSource, Target, EnExList>(state_);
  }

  AState& state_;
};
\end{cppcode}
В конструкторе мы перехватываем по ссылке нераспакованное текущее состояние.
После распаковки мы вызываем функцию
\begin{cppcode}[\nonumbers]
Back::ForceTransition<AState, PSource, Target, EnExList>(state_)
\end{cppcode}
Пока для нас это черный ящик, но надо объяснить, что он делает и какой смысл у шаблонных аргументов.
Данная функция выполняет насильное перемещение машины из состояния \verb"PSource" в состояние \verb"Target".
Она предполагает, что \verb"state_" точно находится в состоянии \verb"PSource".
Данная функция выполняет действие \verb"on_exit" для \verb"PSource" и действие \verb"on_entry" для \verb"Target".
А потому ей нужно знать список \verb"EnExList".
Саму функцию мы обсудим позже.

\paragraph{Визитер \texttt{ProcessOnce}}

И теперь обработчик одного неанонимного перехода:
\begin{cppcode}
template<class Edges, class EnExList, class AState, class Event>
struct ProcessOnce {
  static_assert(areSame2<Event, std::remove_cvref_t<Event>>);
  ProcessOnce(AState& state, Event& event) : state_(state), event_(event) {
  }

  template<class Source>
  bool operator()(Source&) const {
    using PSource = std::remove_cvref_t<Source>;
    return Back::TransitionFrom<AState, PSource, Event, Edges, EnExList>(state_, event_);
  }

  AState& state_;
  Event& event_;
};
\end{cppcode}
Этот обработчик так же перехватывает состояние \verb"state_" в нераспакованном виде, потому что он может его поменять.
А так же он перехватывает ссылку на текущее сообщение \verb"event_".
Внутри же после распаковки вызывается уже знакомая нам функция
\begin{cppcode}[\nonumbers]
Back::TransitionFrom<AState, PSource, Event, Edges, EnExList>(state_, event_);
\end{cppcode}
Которая просто выполняет один переход из \verb"PSource" по сообщению \verb"Event".
Для исполнения нужно знать ребра из \verb"Edges" и действия на вход и выход из \verb"EnExList".

\subsubsection{Имплементация сдвига по выбранному ребру}

\paragraph{Что мы хотим}

Эта функция у нас еще не возникала, но она понадобится для имплементации функций перехода.
Данная функция будет необходима для имплементации других функций перехода.
\begin{cppcode}[\nonumbers]
EdgeTransition<AState, Edge, ArgList, SourceExit, TargetEntry>(state_, ...)
\end{cppcode}
Здесь \verb"Edge" -- ребро по которому надо осуществить переход:
\begin{cppcode}[\nonumbers]
Edge = TL::List<Source, Event, Guard, Action, Target>
\end{cppcode}
Список \verb"ArgList" -- список аргументов не содержащий \verb"Source" и \verb"Target".
В этом списке лежат только неконстантные ссылки.
В этом списке есть ссылка на сообщение, если ребро не анонимное и еще могут лежать ссылки на другие данные, которые используются машиной состояния.
Последние два параметра \verb"SourceExit" и \verb"TargetEntry" -- действия для выхода из \verb"Source" и входа в \verb"Target" соответственно.

\paragraph{Имплементация}

Я следую своему обычному паттерну: служебный шаблон в \verb"detail" с имплементацией и распаковка для удобного доступа к \verb"value":
\begin{cppcode}[\nonumbers]
namespace detail {
template<class AState, class Edge, class ArgList, class SourceExit, class TargetEntry>
struct EdgeTransitionImpl;
} // namespace detail

template<class AState, class Edge, class ArgList, class SourceExit, class TargetEntry>
struct EdgeTransition_t : detail::EdgeTransitionImpl<AState, Edge, ArgList,
                                                     SourceExit, TargetEntry> {
};

template<class AState, class Edge, class ArgList, class SourceExit, class TargetEntry>
constexpr auto EdgeTransition =
    EdgeTransition_t<AState, Edge, ArgList, SourceExit, TargetEntry>::value;
\end{cppcode}
Чтобы суметь имплементировать нужную функцию нам надо распаковать \verb"ArgList".
Это делается вызовом вспомогательного шаблона (это можно было бы сделать и без вспомогательного шаблона специализацией самого шаблона \verb"EdgeTransitionImpl", но я что-то сделал так в два шага).
Вот как выглядит имплементация:
\begin{cppcode}[\nonumbers]
template<class AState, class Edge, class ArgList, class SourceExit, class TargetEntry>
struct EdgeTransitionImpl {
  static constexpr auto value =
      EdgeTransitionUnPack<AState, Edge, ArgList, SourceExit, TargetEntry>::exec;
  using value_type = std::remove_cv_t<decltype(value)>;
};
\end{cppcode}
То есть класс содержит статическую \verb"constexpr" переменную \verb"value", которая является указателем на нужную функцию.
Распаковка устроена так: общая специализация для сообщения об ошибках
\begin{cppcode}
template<class AState, class Edge, class ArgList, class SourceExit, class TargetEntry>
struct EdgeTransitionUnPack {
  static_assert(False<AState>, "Internall Error of EdgeTransition!");
};
\end{cppcode}
и специализация, где происходит распаковка
\begin{cppcode}[\nonumbers]
template<class AState, class Edge, class... Args, template<class...> class List,
         class SourceExit, class TargetEntry>
struct EdgeTransitionUnPack<AState, Edge, List<Args&...>, SourceExit, TargetEntry> {
  /* using declarations */
  
  /* static constants */
  
  static bool exec(AState& current, Args&... args) {
    /*implementation*/
  }
};
\end{cppcode}
В начале идут вспомогательные псевдонимы, чтобы вытащить всю нужную информацию из шаблонных параметров.
В начале вытаскиваем концы ребра перехода и тип сообщение:
\begin{cppcode}[\nonumbers]
using Source = EdgeOp::Source<Edge>;
using Target = EdgeOp::Target<Edge>;
using Event = EdgeOp::Event<Edge>;
\end{cppcode}
Потом достаем \verb"Guard" и \verb"Action" так:
\begin{cppcode}[\nonumbers]
using GuardType = EdgeOp::GuardType<Edge>;
static constexpr auto guard = EdgeOp::Guard<Edge>;
using ActionType = EdgeOp::ActionType<Edge>;
static constexpr auto action = EdgeOp::Action<Edge>;
\end{cppcode}
Теперь вытаскиваем действия на вход и выход
\begin{cppcode}[\nonumbers]
static constexpr auto on_exit = SourceExit::value;
static constexpr auto on_entry = TargetEntry::value;
\end{cppcode}
Напомню, что 
\begin{cppcode}[\nonumbers]
SourceExit  = Value_t<decltype(on_exit), on_exit>;
TargetEntry = Value_t<decltype(on_entry), on_entry>;
\end{cppcode}
и функции доступны по обращению к полю \verb"value".
Для определения корректного поведения нужно еще определить является ли перемещение внутренним (то есть не должно дергать \verb"on_exit" и \verb"on_entry" действия).
Для этого заводим вспомогательную переменную:
\begin{cppcode}[\nonumbers]
static constexpr bool isInternal = EdgeOp::isInternal<Edge>;
\end{cppcode}
Теперь напишем имплементацию для всех случаев:
\begin{cppcode}[\nonumbers]
static bool exec(AState& current, Args&... args) {
  assert(std::get_if<Source>(&current));
  Source& source = *std::get_if<Source>(&current);
  if (!Fun::auto_call(guard, source, args...))
    return false;
  if constexpr (!isInternal) {
    Target target;
    Fun::auto_call(on_exit, source, args...);
    Fun::auto_call(action, source, target, args...);
    Fun::auto_call(on_entry, target, args...);
    current = std::move(target);
  } else {
    Fun::auto_call(action, source, args...);
  }
  return true;
}
\end{cppcode}
В начале независимо от имплементации мы заводим ссылку на распакованный \verb"Source".
Его мы вынимаем из \verb"current" (это куда подставляется \verb"state_".
После мы проверяем условие на \verb"guard".
\begin{cppcode}[\nonumbers]
if (!Fun::auto_call(guard, source, args...))
  return false;
\end{cppcode}
Здесь мы используем шаблон \verb"auto_call", который обсуждался в разделе~\ref{}.
Дело в том, что в этом месте мы переходим к вызову пользовательской функции \verb"guard".
Но мы не знаем в каком порядке и какие аргументы указывал пользователь.
Потому нам нужна автоматическая диспетчеризация всех аргументов, которую и решает \verb"auto_call".
Если мы используем данные из State Machine, которые запакованы в \verb"Tuple", то надо еще дополнительно добавить распаковку элементов \verb"Tuple" в отдельные аргументы.
То есть логически данный вызов эквивалентен:
\begin{cppcode}[\nonumbers]
if (!guard(source, args...))
  return false;
\end{cppcode}
Но только с правильным порядком аргументов (с выкинутыми ненужными).
Если \verb"guard" не сработал, то мы тут же выходим с сообщением \verb"false".
А дальнейшее поведение зависит от того какой у нас вид перехода:
\begin{enumerate}
\item Первый случай -- не внутренний переход.
Тут надо сделать все необходимые шаги.
\begin{enumerate}
\item Создать новое состояние
\begin{cppcode}[\nonumbers]
Target target;
\end{cppcode}

\item Выполнить действия на выход для текущего состояния.
Здесь так же нужна диспетчеризация:
\begin{cppcode}[\nonumbers]
Fun::auto_call(on_exit, source, args...);
\end{cppcode}

\item Выполнить действие перехода
\begin{cppcode}[\nonumbers]
Fun::auto_call(action, source, target, args...);
\end{cppcode}

\item Выполнить действия на вход в новое состояние
\begin{cppcode}[\nonumbers]
Fun::auto_call(on_entry, target, args...);
\end{cppcode}

\item Заменить старое состояние на новое
\begin{cppcode}[\nonumbers]
 current = std::move(target);
\end{cppcode}
и вернуть \verb"true", так как переход произошел.
\end{enumerate}

\item Второй случай -- внутренний переход, который не активирует действия на вход и выход.
А потому у нас выполняется только \verb"action":
\begin{cppcode}[\nonumbers]
Fun::auto_call(action, source, args...);
\end{cppcode}
и возвращаем \verb"true".
\end{enumerate}

\subsubsection{Имплементация сдвига по неизвестному ребру}

\paragraph{Что мы хотим}

Теперь надо имплементировать функцию:
\begin{cppcode}[\nonumbers]
Back::TransitionFrom<AState, Source, Event, Edges, EnExList>(state_, event_, ...);
\end{cppcode}
В начале несколько замечаний о работе функции.
Функция предполагает, что внутри \verb"state_" находится состояние \verb"Source".
Переменная \verb"event" имеет тип \verb"Event".
Список \verb"Edges" содержит все возможные ребра, а данная функция выбирает все ребра из него, которые начинаются с \verb"Source" и имеют сообщение типа \verb"Event".
После чего перебором по всем ребрам проверяется \verb"Guard".
И производится переход на первый сработавший \verb"Guard".
Дополнительные аргументы в \verb"..." означают возможное наличие дополнительных данных, которые лежат в State Machine.
Если внутри машины состояний данные лежат в \verb"Tuple" с именем \verb"data_", то вызов в общем случае будет
\begin{cppcode}[\nonumbers]
Back::TransitionFrom<AState, Source, Event, Edges, EnExList>(state_, event_, data_);
\end{cppcode}
и уже эта функция должна распаковать \verb"data_" и подставить что куда надо автоматически.
В текущей имплементации я буду просто игнорировать дополнительные аргументы, но код будет написан так, что его легко будет расширить до большего списка аргументов.
Эта функция всегда принимает нераспакованное \verb"state" в качестве первого аргумента.
Порядок остальных аргументов определяется автоматически.
Все аргументы передаются по неконстантной ссылке.

\paragraph{Имплементация}

Имплементация опять делается в два шага: служебный шаблон с имплемнтацией и удобный шаблон для обращения к функции:
\begin{cppcode}
namespace detail {
template<class AState, class Source, class Event, class Edges, class AStatesEnEx>
struct TransitionFromImpl;
} // namespace detail

template<class AState, class Source, class Event, class Edges, class StatesEnEx>
constexpr auto TransitionFrom =
    detail::TransitionFromImpl<AState, Source, Event, Edges, StatesEnEx>::value;
\end{cppcode}
Сама имплементация устроена так:
\begin{cppcode}
template<class AState, class Source, class Event, class Edges, class StatesEnEx>
struct TransitionFromImpl {
  using EdgesOut = EdgesFromBy<Edges, Source, Event>;
  using States = TL::NoDuplicates<
      TL::PushFront1<Source, ExtractStates<EdgesOut, TL::EmptyLike<EdgesOut>>>>;

  using EnExFromList =
      TL::Filter<StateOp::Bind<States>::template Contains_t, StatesEnEx>;

  using Args = ExtractAllArgs<EdgesOut, EnExFromList>;
  using Refs = MakeNonConstLRef<Args>;
  using RefsNoStates = TL::Diff<Refs, MakeNonConstLRef<States>>;
  static constexpr bool isAnonymous = areSame2<Event, void>;

  using ArgList = IfElse_u<isAnonymous, Type_t<RefsNoStates>,
                           TL::NoDuplicates_w<TL::PushFront1_w<
                               std::add_lvalue_reference<Event>, Type_t<RefsNoStates>>>>;
  static constexpr auto value =
      TransitionFromUnPack<AState, Source, ArgList, EdgesOut, EnExFromList>::exec;
  using value_type = std::remove_cv_t<decltype(value)>;
};
\end{cppcode}
Давайте пройдемся по имплементации:
\begin{enumerate}
\item Начнем с главного в самом конце:
\begin{cppcode}[\nonumbers]
static constexpr auto value =
    TransitionFromUnPack<AState, Source, ArgList, EdgesOut, EnExFromList>::exec;
using value_type = std::remove_cv_t<decltype(value)>;
\end{cppcode}
Мы берем функцию \verb"exec" из вспомогательного шаблона, который просто распаковывает \verb"ArgList".
И еще запоминаем ее тип в \verb"value_type".
Но перед вызовом этого шаблона мы должны сделать подготовительную работу: надо найти какие аргументы понадобятся внутри при переходе по каждому ребру, которое исходит из \verb"Source".
Для этого и нужны псевдонимы.

\item В начале отбираем все ребра выходящие в точности из \verb"Source" и имеющие тип сообщения \verb"Event".
Я это делаю с помощью вспомогательного шаблона, который еще не обсуждался (его имплементация будет ниже):
\begin{cppcode}[\nonumbers]
using EdgesOut = EdgesFromBy<Edges, Source, Event>;
\end{cppcode}

\item Потом достаем все вершины, которые затронуты ребрами из \verb"EngesOut" (сюда входит \verb"Source" и все \verb"Target" по всем ребрам).
\begin{cppcode}[\nonumbers]
using States = TL::NoDuplicates<
    TL::PushFront1<Source, ExtractStates<EdgesOut, TL::EmptyLike<EdgesOut>>>>;
\end{cppcode}
Здесь надо быть аккуратным.
Так как если множество ребер пусто, то мы извлекаем пустое множество вершин.
В частности \verb"Source" надо в этом случае добавить искусственно и потом убрать дубликаты.

\item Теперь надо оставить только те записи \verb"EnExList", в которых вершины из списка \verb"States".
\begin{cppcode}[\nonumbers]
using EnExFromList =
    TL::Filter<StateOp::Bind<States>::template Contains_t, StatesEnEx>;
\end{cppcode} 

\item Теперь надо вычленить все возможные аргументы из всех возможных действий, которые у нас могут произойти.
То есть надо взять все \verb"guard"-ы, \verb"action"-ы, \verb"on_exit" для \verb"Source" и \verb"on_entry" для каждого \verb"Target" и из сигнатуры всех этих функций выбрать типы всех аргументов.
Я это делаю отдельным шаблоном, который определю позже
\begin{cppcode}[\nonumbers]
using Args = ExtractAllArgs<EdgesOut, EnExFromList>;
\end{cppcode} 

\item Сейчас аргументы могут быть разного вида.
Для каждого типа \verb"T" аргумент может встретиться как \verb"T" или \verb"T&" или \verb"const T&" или \verb"const T".
Я заменяю все аргументы на неконстантную ссылку:
\begin{cppcode}[\nonumbers]
using Refs = MakeNonConstLRef<Args>;
\end{cppcode}
Это разумно, потому что в самом верху стека вызова у меня в качестве аргументов будут выступать именно lvalue.
Это либо \verb"state_", либо созданная вершина \verb"target", либо \verb"event" являющийся аргументом функции \verb"process", либо данные \verb"data_", которые лежат внутри машины состояний.

\item После чего я удаляю все состояния из общего списка аргументов.
Это связано с тем, чтобы избежать дублирование типов, если \verb"Source" и \verb"Target" одно и то же.
Кроме того, \verb"Source" всегда вычленяется из \verb"AState".
Так что с состояниями я буду разбираться отдельно.
\begin{cppcode}[\nonumbers]
using RefsNoStates = TL::Diff<Refs, MakeNonConstLRef<States>>;
\end{cppcode}
Так как в \verb"Refs" в находятся неконстантные ссылки без повторений, то надо просто вычесть список ссылок на состояния.

\item  Теперь мне надо понять нужно ли передавать \verb"event" как аргумент или нет.
Это проверяется флагом:
\begin{cppcode}[\nonumbers]
static constexpr bool isAnonymous = areSame2<Event, void>;
\end{cppcode} 
Без него у меня будет \verb"void&" добавлен в \verb"ArgList", что конечно же будет ошибка компиляции.

\item Теперь добавим ко всем аргументам \verb"Event&", если у нас не анонимный переход и избавимся от дубликатов:
\begin{cppcode}[\nonumbers]
using ArgList =
    IfElse_u<isAnonymous, Type_t<RefsNoStates>,
             TL::NoDuplicates_w<TL::PushFront1_w<
                 std::add_lvalue_reference<Event>, Type_t<RefsNoStates>>>>;
\end{cppcode}
Мне приходится тут производить вычисления внутри \verb"IfElse" не распаковывая типы, иначе было бы сообщение об ошибке.
Более подробно, я бы мог написать так:
\begin{cppcode}[\nonumbers]
using ArgList =
    IfElse<isAnonymous,
           RefsNoStates,
           TL::NoDuplicates<TL::PushFront1<std::add_lvalue_reference_t<Event>,
                                           RefsNoStates>>>;
\end{cppcode}
Компилятор в этом случае должен вычислить все аргументы.
В частности он должен вычислить 
\begin{cppcode}[\nonumbers]
std::add_lvalue_reference_t<Event>
\end{cppcode}
Но тогда бы в случае анонимного сообщения \verb"Event" было бы \verb"void", то компилятор бы вычислял
\begin{cppcode}[\nonumbers]
std::add_lvalue_reference_t<void>
\end{cppcode}
что привело бы к ошибке компиляции.
Для того, чтобы избежать этой проблемы, надо запрятать результат вычислений внутрь другого шаблона каковым как раз является
\begin{cppcode}[\nonumbers]
std::add_lvalue_reference<void>
\end{cppcode}
И вот это не приведет к ошибке компиляции.
Но тогда снаружи надо использовать операции над запакованным типом, которые у меня обозначающся с суффиксом \verb"_w".
А распаковка делается с помощью \verb"_u" версии (как раз в \verb"IfElse" этот суффикс).
\end{enumerate}
Теперь надо распаковать \verb"ArgList", что делается вспомогательным шаблоном.
Для экономии времени и места, я напишу лишь содержательную часть, без пробверки ошибок:
\begin{cppcode}
template<class AState, class Source, template<class...> class TList,
         class... Args, class EdgesOut, class EnExList>
struct TransitionFromUnPack<AState, Source, TList<Args&...>, EdgesOut,
                            EnExList> {
  /* using declarations to define FList */

  static bool exec(AState& current, Args&... args) {
    return Fun::cascade_call<FList>(current, args...) != -1;
  }
};
\end{cppcode}
Я специально скрыл вспомогательные шаблоны, чтобы лучше было видно, что тут происходит.
Функция \verb"exec", которая должна выполнить переход по ребру из списка, вызывает \verb"cascade_call" (см.~\ref{}).
Эта функция принимает список функций \verb"FList" для исполнения выполняет функции строго по списку, подставляя в них переданные аргументы (данная версия без диспетчеризации).
Как только получен ответ \verb"true", функция \verb"cascade_call" завершает исполнение и возвращает номер элемента в списке, который выдал \verb"true".
Если никто не выдал \verb"true", то возвращается значение $-1$.
Так как \verb"exec" должна вернуть флаг был ли переход, то мы просто сравниваем с $-1$ на неравенство.
Теперь нам надо построить список \verb"FList", в котором лежат функции, выполняющие переход вдоль одного ребра, для всех ребер из \verb"EdgesOut".
Построением такого списка и занимаются \verb"using" директивы в начале класса.
В этом месте нам понадобится шаблон шаблон \verb"EdgeTransition" разобранный в предыдущем разделе.

Давайте посмотрим как работает построение нужного списка действий:
\begin{enumerate}
\item В начале я завожу псевдоним для списка переданных аргументов
\begin{cppcode}[\nonumbers]
using ArgList = TList<Args&...>;
\end{cppcode}
Просто для удобства.

\item Я хочу для каждого ребра взять его начало и у начала взять 

\item Теперь вытаскиваем все записи из \verb"EnExList", в которых встречается \verb"Source" ребра.
\begin{cppcode}[\nonumbers]
using EdgeSourceEnExList =
    TL::Filter<StateOp::Bind<Source>::template isSame_t, EnExList>;
\end{cppcode}
Получим список из в точности $1$ элемента.

\item Теперь вытаскиваем \verb"on_exit" для \verb"Source":
\begin{cppcode}[\nonumbers]
using SourceExit = StateOp::OnExitW<TL::Front<EdgeSourceEnExList>>;
\end{cppcode}

\item Концы у ребер уже все различные.
Потому для каждого ребра надо найти свою запись в \verb"EnExList" и вытащить его \verb"on_entry" действие.
Потому далее я все операции делаю зависимыми от ребра.

\item Вот такая процедура вытаскивает для каждого ребра список из одной записи:
\begin{cppcode}[\nonumbers]
template<class Edge>
using EdgeTargetEnExList =
    TL::Filter<StateOp::Bind<EdgeOp::Target<Edge>>::template isSame_t, EnExList>;
\end{cppcode}
Тут надо сказать, что \verb"EdgeOp::Target" должна не просто брать $4$-ый элемент, а она проверяет, если $4$-ый элемент \verb"void", то она берет \verb"Source", иначе $4$-ый элемент ребра.

\item Вот такая процедура из полученного списка берет единственный элемент и извлекает \verb"on_entry" действие.
\begin{cppcode}[\nonumbers]
template<class Edge>
using TargetEntry = StateOp::OnEntryW<TL::Front<EdgeTargetEnExList<Edge>>>;
\end{cppcode}

\item Теперь вспомогательный шаблон, который для каждого ребра строит функцию перехода по этому ребру.
А мы уже обсуждали в точности эту задачу.
Для ее решения используется шаблон \verb"EdgeTransition_t", куда мы подставляем нужные данные:
\begin{cppcode}[\nonumbers]
template<class Edge>
using EdgeTransitionT = EdgeTransition_t<AState, Edge, ArgList,
                                         SourceExit, TargetEntry<Edge>>;
\end{cppcode}

\item Выше у нас получился список из типов, в которые завернуты функции.
Теперь надо из списка типов сделать список функций перекодированием:
\begin{cppcode}[\nonumbers]
using FList = AL::FromTList<FListT>;
\end{cppcode}
И на этом подготовка к вызову закончена.
\end{enumerate}

\subsubsection{Имплементация насильного перемещения}

Теперь надо имплементировать функцию
\begin{cppcode}[\nonumbers]
Back::ForceTransition<AState, Source, Target, EnExList>(state_, ...);
\end{cppcode}
Эта функция считает что \verb"state_" находится в состоянии \verb"Source".
Она не проверяет это условие.
Обязательные аргументы -- \verb"state_" для нераспакованного \verb"Source".
Дополнительные аргументы отмеченные \verb"..." могут отвечать за данные State Machine, которые указаны в аргументах действий \verb"on_exit" для \verb"Source" и \verb"on_entry" для \verb"Target".

\paragraph{Имплементация}

Имлементация опять делается в два шага: служебный шаблон с имплементацией и удобный шаблон для обращения к функции:
\begin{cppcode}
namespace detail {
template<class AState, class Source, class Target, class EnExList>
struct ForceTransitionImpl;
} // namespace detail

template<class AState, class Source, class Target, class EnExList>
constexpr auto ForceTransition =
    detail::ForceTransitionImpl<AState, Source, Target, EnExList>::value;

\end{cppcode}
Сама имплементация выглядит так:
\begin{cppcode}

template<class AState, class Source, class Target, class EnExList>
struct ForceTransitionImpl {
  /* using declarations */
  /* define SourceNode, TargetNode, RefsNoStates */

  static constexpr auto value =
      ForceTransitionUnPack<AState, Source, Target, SourceNode, TargetNode,
                            RefsNoStates>::exec;

  using value_type = std::remove_cv_t<decltype(value)>;
};
\end{cppcode}
В начале шаблона строятся три вспомогательных класса:
\begin{enumerate}
\item \verb"SourceNode" -- содержит \verb"on_entry" и \verb"on_exit" методы для \verb"Source".

\item \verb"TargetNode" -- содержит \verb"on_entry" и \verb"on_exit" методы для \verb"Target".

\item \verb"RefsNoStates" -- список аргументов для \verb"exec" исключающий состояния \verb"Source" и \verb"Target".
\end{enumerate}
Давайте напишем, как происходят эти вычисления:
\begin{enumerate}
\item В начале надо вытащить единственную запись для вершины \verb"Source":
\begin{cppcode}[\nonumbers]
using SourceNode =
    TL::Front<TL::Filter<StateOp::Bind<Source>::template isSame_t, EnExList>>;
\end{cppcode}

\item Аналогично для вершины \verb"Target":
\begin{cppcode}[\nonumbers]
using TargetNode =
    TL::Front<TL::Filter<StateOp::Bind<Target>::template isSame_t, EnExList>>;
\end{cppcode}

\item  Теперь извлекаем из вершины для \verb"Source" тип действия \verb"on_exit" и из него берем аргументы с помощью шаблона \verb"ArgsOf", который мы разбирали в разделе~\ref{section::ArgsOf}.
Аналогично поступаем с вершиной \verb"Target", но берем \verb"on_entry" действие.
А потом оба списка аргументов сливаем в один и убираем дубликаты.
\begin{cppcode}[\nonumbers]
using ExitArgs = Fun::ArgsOf<StateOp::OnExitType<SourceNode>>;
using EntryArgs = Fun::ArgsOf<StateOp::OnEntryType<TargetNode>>;
using Args = TL::NoDuplicates<TL::Merge2<ExitArgs, EntryArgs>>;
\end{cppcode}

\item Теперь трансформируем все аргументы в неконстантные ссылки (эта функция еще и убирает дубликаты):
\begin{cppcode}[\nonumbers]
using Refs = MakeNonConstLRef<Args>;
\end{cppcode}

\item И теперь удаляем ссылки на \verb"Source" и \verb"Target".
\begin{cppcode}[\nonumbers]
using RefsNoStates = TL::Diff<Refs, TL::List<Source&, Target&>>;
\end{cppcode}
\end{enumerate}
Теперь напишем как выглядит шаблон с распаковкой \verb"RefsNoStates" определяющий \verb"exec".
Я сразу перейду к содержательной специализации.
\begin{cppcode}
template<class AState, class Source, class Target, class SourceNode,
         class TargetNode, class... Args, template<class...> class List>
struct ForceTransitionUnPack<AState, Source, Target, SourceNode, TargetNode,
                             List<Args&...>> {
  static constexpr auto on_exit = StateOp::OnExit<SourceNode>;
  static constexpr auto on_entry = StateOp::OnEntry<TargetNode>;

  static bool exec(AState& current, Args&... args) {
    assert(std::get_if<Source>(&current));
    Source& source = *std::get_if<Source>(&current);
    Target target;
    Fun::auto_call(on_exit, source, args...);
    Fun::auto_call(on_entry, target, args...);
    current = std::move(target);
    return true;
  }
};
\end{cppcode}
Тут все прямолинейно.
Мы в начале вынимаем сами действия \verb"on_exit" и \verb"on_entry" из шаблонных параметров.
А в самом методе создаем новое состояние, потом выход из \verb"source", вход в \verb"target" и делаем подмену внутри \verb"std::variant".
Тут наверное можно было бы рассмотреть случай, когда \verb"Source" и \verb"Target" одно и то же и избежать создания нового состояния \verb"target".
Но если после работы \verb"on_exit" состояние совпадает с состоянием после дефолтного конструктора, то это не требуется.
Но чтобы было наверняка, мы просто заводим новое состояние с дефолтным значением и все.


\subsubsection{Вспомогательные методы}

Давайте я для полноты картины опишу некоторые вспомогательные шаблоны, которые были использованы выше.

\paragraph{Отсеивание ребер}

Напомню, что задача заключается в следующем.
Есть список ребер
\begin{cppcode}[\nonumbers]
using Edges = TL::List<Edge1, Edge2, ...>;
\end{cppcode}
Теперь есть три задачи:
\begin{enumerate}
\item Отобрать только те ребра, что ведут из вершины \verb"State".

\item Отобрать только те ребра, что реагируют на сообщения типа \verb"Event".

\item Отобрать ребра удовлетворяющие обоим условиям.
\end{enumerate}
Понятно, что тут надо просто использовать \verb"TL::Filter" с нужным предикатом.
Вот как выглядит имплементация отборки по \verb"Source":
\begin{cppcode}[\nonumbers]
namespace detail {
template<class TList, class State, bool isCorrect = isEdgeList<TList>>
struct EdgesFromImpl {
  static_assert(isEdgeList<TList>, "The first argument of EdgesFrom MUST be an edge list!");
};

template<class TList, class State>
struct EdgesFromImpl<TList, State, true> {
  using type = TL::Filter<EdgeOp::Bind<State>::template isSourceSame_t, TList>;
};
} // namespace detail

template<class TList, class State>
struct EdgesFrom_t : detail::EdgesFromImpl<std::remove_cv_t<TList>,
                                           std::remove_cvref_t<State>> {};

template<class TList, class State>
using EdgesFrom = typename EdgesFrom_t<TList, State>::type;
\end{cppcode}
Абсолютно аналогично имплементируется отборка по сообщению:
\begin{cppcode}[\nonumbers]
namespace detail {
template<class TList, class Event, bool isCorrect = isEdgeList<TList>>
struct EdgesByImpl {
  static_assert(isEdgeList<TList>, "The first argument of EdgesBy MUST be an edge list!");
};

template<class TList, class Event>
struct EdgesByImpl<TList, Event, true> {
  using type = TL::Filter<EdgeOp::Bind<Event>::template isEventSame_t, TList>;
};
} // namespace detail

template<class TList, class Event>
struct EdgesBy_t
    : detail::EdgesByImpl<std::remove_cv_t<TList>, std::remove_cvref_t<Event>> {
};

template<class TList, class Event>
using EdgesBy = typename EdgesBy_t<TList, Event>::type;
\end{cppcode}
И теперь можно сразу имплементировать последний метод как композицию предыдущих двух даже без функции имплементации.%
\footnote{По хорошему все равно надо бы сделать имплементацию, чтобы сообщать об ошибках относительно имени этой функции.
Сейчас же пользователь увидит ошибки от функций \texttt{EdgesFrom} и \texttt{EdgesBy} и будет сложно понять, что это именно функция \texttt{EdgesFromBy}.
Но я как обычно ленюсь.}
\begin{cppcode}[\nonumbers]
template<class TList, class State, class Event>
struct EdgesFromBy_t : EdgesBy_t<EdgesFrom<TList, State>, Event> {};

template<class TList, class State, class Event>
using EdgesFromBy = typename EdgesFromBy_t<TList, State, Event>::type;
\end{cppcode}

\paragraph{Вычленение зависимостей}

Напомню, что вся машина состояний задается двумя списками:
\begin{cppcode}[\nonumbers]
using Edges = TL::List<Edge1, Edge2, ...>;
using EnExList = TL::List<Node1, Node2, ...>;
\end{cppcode}
Первые содержат \verb"Guard" и \verb"Action", а вторые \verb"OnExit" и \verb"OnEntry" действия.
Все их аргументы являются зависимостями, которые надо автоматически подтягивать.
Мы просто вытянем списки типов для всех этих действий, на	йдем их аргументы и сольем все списки вместе.
Потом уберем все повторения.
Вот так выглядит имплементация:%
\footnote{Да, я опять забил на проверку ошибок.}
\begin{cppcode}[\nonumbers]
template<class Edges, class StatesEnEx>
struct ExtractAllArgs_t {
  using GuardTypeList = ExtractGuardTypes<Edges>;
  using ActionTypeList = ExtractActionTypes<Edges>;
  using OnEntryTypeList = ExtractOnEntryTypes<StatesEnEx>;
  using OnExitTypeList = ExtractOnExitTypes<StatesEnEx>;
  using FList = TL::Merge<GuardTypeList, ActionTypeList, OnEntryTypeList, OnExitTypeList>;
  using type = detail::ExtractArgs<FList>;
};
\end{cppcode}
Я не буду определять все четыре функции вычленения типов для действий.
Обойдусь первой.
\begin{cppcode}[\nonumbers]
template<class Edges>
struct ExtractGuardTypes_t
    : TL::NoDuplicates_t<TL::Map<EdgeOp::GuardType, Edges>> {};

template<class Edges>
using ExtractGuardTypes = typename ExtractGuardTypes_t<Edges>::type;
\end{cppcode}
Аналогично пишутся остальные методы, учитывая, что у нас есть уже все вспомогательные методы.
Теперь метод по вычленению аргументов из списка сигнатур функций:
\begin{cppcode}[\nonumbers]
template<class FList>
struct ExtractArgs_t : NoDuplicates_t<TL::Join<TL::Map<Fun::ArgsOf, FList>>> {
};

template<class FList>
using ExtractArgs = typename ExtractArgs_t<FList>::type;
\end{cppcode}
Тут надо сказать про метод \verb"TL::Join", который я не обсуждал.
Этот метод ожидает, что на вход ему подают список списков:
\begin{cppcode}[\nonumbers]
using List1 = List<List<A, B, C>, List<>, List<D, E>>;
\end{cppcode}
И он сливает их в один большой список.
\begin{cppcode}[\nonumbers]
using List2 = TL::Join<List1>;
// List2 = List<A, B, C, D, E>;
\end{cppcode}
Я надеюсь, что этот метод вы тоже догадаетесь как имплементировать.
Он как минимум сводится к \verb"TL::Merge".

\paragraph{Нормализация типов}

Этот метод пишется максимально в лоб.
Сначала снимаем со всех аргументов ссылки и \verb"const/volatile" квалификаторы, а потом надеваем lvalue ссылки и убираем дубликаты.
Я опять поленился с проверкой.
\begin{cppcode}[\nonumbers]
template<class TList>
struct MakeNonConstLRef_t {
  using NoCVRefList = TL::Map<std::remove_cvref_t, TList>;
  using RefList = TL::Map<std::add_lvalue_reference_t, NoCVRefList>;
  using type = TL::NoDuplicates<RefList>;
};

template<class TList>
using MakeNonConstLRef = typename MakeNonConstLRef_t<TList>::type;
\end{cppcode}

\subsection{Проверка таблицы на корректность}

Так как наша машина состояний задается полностью во время компиляции, то мы можем сделать почти все проверки корректности ее построения во время компиляции.
При этом мы уже сделали некоторые проверки на уровне FrontEnd-а, когда сообщали пользователю о недопустимых операциях.
Однако, на уровне FrontEnd-а не хочется делать более серьезные проверки не относящиеся к DSL.
С другой стороны при построении машины состояний мы уже хотим иметь корректные данные.
Давайте вспомним какие данные принимает машина состояний после конвертации данных из FrontEnd-а:
\begin{cppcode}[\nonumbers]
template<class Edges, class EnExList, class States, class InitStates,
         class FinalStates, class Events, bool isCorrect = /*checks*/>
class SM {
  /* Implementation */
};
\end{cppcode}
Теперь я хочу перечислить условия, которые у нас точно гарантируются методами конвертации:
\begin{enumerate}
\item Корректность \verb"States".
Гарантируется, что список всех состояний встречающихся в \verb"Edges" и \verb"EnExList" совпадает со \verb"States".
Кроме того, все типы в этом списке не ссылки и без \verb"const/volatile" квалификаторов.

\item Корректность \verb"Events".
Это значит, что список всех типов сообщений, которые встречаются в \verb"Edges", в точности совпадает с \verb"Events".
Кроме того, все типы в этом списке не ссылки и без \verb"const/volatile" квалификаторов.

\item Список \verb"EnExList" содержит в точности по одной записи для каждого состояния \verb"States".
Вот эту проверку можно было бы убрать из конвертирующего кода.
Она подразумевается в момент слияния вершин списка для одного и того же состояния (см.~\ref{section::SMFrontToBackConvert}).

\item Так же мы гарантируем, что \verb"InitialStates" и \verb"FinalStates" является подмножествами в \verb"States".
\end{enumerate}
Что мы хотим и можем проверить во время компиляции:
\begin{enumerate}
\item Список \verb"InitialStates" содержит в точности одну вершину.

\item Список \verb"InitialStates" и \verb"FinalStates" не пересекаются.

\item Списки \verb"States", \verb"Events" не пересекаются.

\item Для каждого состояния \verb"S" из списка \verb"States" либо есть только одно анонимное ребро ведущее из него, либо  у всех анонимных ребер нетривиальные \verb"Guard"-ы.
Это не дает полных гарантий, поскольку сравнивать \verb"Guard"-ы мы можем только по указателям на функции и пользователь мог задать что-то плохое, что можно отловить только в run-time.
Но как минимум условия выше проверяются во время компиляции.

\item Если в графе переходов оставить только анонимные ребра, то в полученном графе нет петель.
Это необходимо, потому что переход по анонимным ребрам выполняется автоматически и наличие таких петель автоматически уведет машину в бесконечный цикл.
\end{enumerate}
Условия которые мы не можем проверить во время компиляции.
\begin{enumerate}
\item По-хорошему для машины состояний требуется, чтобы для всех ребер выходящих из фиксированной вершины не возможно одновременное срабатывание двух \verb"Guard"-ов, то есть все переходы взаимоисключающие.
Такое просто так отследить не получится, особенно если у вас состояния и сообщения содержат run-time информацию.
Но если очень хочется, то можно поставить \verb"assert"-ы на это условие.
Это потребует переработки имплементации методов перехода по ребрам, но это не сложно.
\end{enumerate}
Я не буду писать тривиальные проверки, а напишу только две интересные.

\subsubsection{Корректность анонимных ребер}

В начале надо проверить условие для отдельной вершины \verb"S", что если из нее выходит более одного анонимного ребра, то у них у всех нетривиальные гарды.
Как обычно двухшаговая имплементация.
В начале сообщения об ошибках в общем шаблоне:%
\footnote{Тут я использую функцию для проверки, что нам действительно дан список ребер.
Ее я обсужу ниже.}
\begin{cppcode}[\nonumbers]
template<class Edges, class F, bool isCorrect = isEdgeList<Edges>>
struct areAnonymousValidAtImpl {
  static_assert(isEdgeList<Edges>,
                "The first argument of areAnonymousValidAt MUST be a list of edges!");
};
\end{cppcode}
И теперь сама имплементация:
\begin{cppcode}[\nonumbers]
template<class Edges, class F>
struct areAnonymousValidAtImpl<Edges, F, true> {
  using AnonymousEdgesAt = EdgesFromBy<Edges, F, void>;
  template<class Edge>
  struct hasNonTrivialGuard_t : Not_w<areSame2_t<EdgeOp::GuardW<Edge>, NoGuard>> {};
  
  using value_type = bool;
  static constexpr bool value =
      Imply<(TL::Size<AnonymousEdgesAt> > 1), TL::ifAll<hasNonTrivialGuard_t, AnonymousEdgesAt>>;
};
\end{cppcode}
Анонимные ребра я вычленяю по типу события, в этом случае \verb"Event" будет \verb"void".
Тут обратите внимание, как я сравниваю гарды.
Я вынимаю их обертки
\begin{cppcode}[\nonumbers]
Value_t<decltype(guard), guard>
\end{cppcode}
и сравниваю их на равенство.%
\footnote{Если вы будете просто сравнивать указатели на функции на равенство, то у вас есть несколько проблем.
Во-первых, указатели на функции разных типов сравнивать на равенство нельзя и вам ругнется компилятор.
Во-вторых, у вас может быть не функция, а объект с перегруженным оператором скобки.
Тогда вообще проблема сравнивать два таких адреса.}
Для тривиального действия я использую заранее подготовленный псевдоним для обертки \verb"NoGuard".
И теперь как обычно публичная часть для использования:
\begin{cppcode}[\nonumbers]
template<class Edges, class F>
struct areAnonymousValidAt_t : detail::areAnonymousValidAtImpl<Edges, F> {};

template<class Edges, class F>
constexpr bool areAnonymousValidAt = areAnonymousValidAt_t<Edges, F>::value;
\end{cppcode}
Теперь надо сделать проверку корректности для всех вершин.
\begin{cppcode}[\nonumbers]
template<class Edges, bool isCorrect = isEdgeList<Edges>>
struct areAnonymousValidImpl {
  static_assert(isEdgeList<Edges>,
                "The first argument of areAnonymousValid MUST be a list of edges!");
};

template<class Edges>
struct areAnonymousValidImpl<Edges, true> {
  using States = ExtractStates<Edges, TL::List<>>;
  template<class S>
  using hasValidEdgesAt = areAnonymousValidAt_t<Edges, S>;

  using value_type = bool;
  static constexpr value_type value = TL::ifAll<hasValidEdgesAt, States>;
};
\end{cppcode}
Тут я не мудрствую лукаво, я просто вычленяю все вершины из списка ребер считая, что \verb"EnExList" пустой.
А дальше просто проверяю, что условие выполнено для всех вершин.
И конечно же публичная часть:
\begin{cppcode}[\nonumbers]
template<class Edges>
struct areAnonymousValid_t : detail::areAnonymousValidImpl<Edges> {};

template<class Edges>
constexpr bool areAnonymousValid = areAnonymousValid_t<Edges>::value;
\end{cppcode}

\subsubsection{Отсутствие петель в анонимных переходах}

Я не буду особенно оптимизировать код проверки.
Я сделаю обычный DFS для каждой вершины графа.
То есть я буду стартовать в конкретной вершине \verb"S" и потом идти в глубину и запоминать предыдущий путь, чтобы понять вернулись ли мы в одну из пройденных вершин или нет.
Вот как выглядит код имплементации:
\begin{cppcode}[\nonumbers]
template<class Edges, class S, class Visited = TL::EmptyLike<Edges>>
struct hasNoCycleFrom_t {
  static constexpr bool isSVisited = TL::Contains<Visited, S>;

  template<class S1, class Visited1>
  struct hasNoCycleFromNonVis_t {
    using EdgesFromS = EdgesFrom<Edges, S1>;
    using Visited2 = TL::Append<Visited1, S1>;
    template<class Edge>
    using hasNoCycleAfter_t = hasNoCycleFrom_t<Edges, EdgeOp::Target<Edge>, Visited2>;
    
    using value_type = bool;
    static constexpr value_type value = TL::ifAll<hasNoCycleAfter_t, EdgesFromS>;
  };

  using value_type = bool;
  static constexpr value_type value =
      IfElse_v<isSVisited, Bool_t<false>, hasNoCycleFromNonVis_t<S, Visited>>;
};
\end{cppcode}
Здесь \verb"Visited" -- вершины составляющие пройденный путь до данной вершины \verb"S".
В начале этот список пустой.
Дальше я проверяю вершина \verb"S" была посещена или нет.
Если да, то мы нашли цикл и надо вернуть \verb"false".
Если же вершина посещена не была, то я добавляю ее в \verb"Visited" и запускаю рекурсивно проверку для каждого исходящего ребра.
Проверка для непосещенной вершины вынесена в отдельный шаблон:
\begin{cppcode}[\nonumbers]
template<class S1, class Visited1>
struct hasNoCycleFromNonVis_t {
  using EdgesFromS = EdgesFrom<Edges, S1>;
  using Visited2 = TL::Append<Visited1, S1>;
  template<class Edge>
  using hasNoCycleAfter_t = hasNoCycleFrom_t<Edges, EdgeOp::Target<Edge>, Visited2>;
  
  using value_type = bool;
  static constexpr value_type value = TL::ifAll<hasNoCycleAfter_t, EdgesFromS>;
};
\end{cppcode}
Тут в точности все как я описал выше.
После этого я завожу непосредственно переменную с результатом (не факт, что она мне будет нужна):
\begin{cppcode}[\nonumbers]
template<class Edges, class S>
constexpr bool hasNoCycleFrom = hasNoCycleFrom_t<Edges, S>::value;
\end{cppcode}
Теперь имплементация проверки, что граф является Directed Acyclic Graph.
\begin{cppcode}[\nonumbers]
template<class Edges, bool isCorrect = isEdgeList<Edges>>
struct isDAGImpl {
  static_assert(isEdgeList<Edges>, "The argument of isDAG MUST be a list of edges!");
};

template<class Edges>
struct isDAGImpl<Edges, true> {
  using States = ExtractStates<Edges, TL::EmptyLike<Edges>>;
  template<class S>
  using hasNoCycleFromT = hasNoCycleFrom_t<Edges, S>;

  using value_type = bool;
  static constexpr value_type value = TL::ifAll<hasNoCycleFromT, States>;
};
\end{cppcode}
Тут без сюрпризов, я вычленяю все состояния из ребер и после этого для каждого состояния запускаю проверку, что из него нет цикла, и проверяю, что все условия выполнены одновременно.
Думаю можно наловчиться и сделать еще помеченные состояния, в которых мы уже были, чтобы не перебирать лишние состояния, но я пожалуй тут не буду этим заниматься.
И как обычно методы для использования:
\begin{cppcode}[\nonumbers]
template<class Edges>
struct isDAG_t : detail::isDAGImpl<Edges> {};

template<class Edges>
constexpr bool isDAG = isDAG_t<Edges>::value;
\end{cppcode}

\subsubsection{Проверка корректности ребер}

% TO DO
TO DO
\begin{cppcode}[\nonumbers]

\end{cppcode}


\subsection{Примеры использования}

% TO DO
TO DO

}

\end{document}
